% !TEX program = xelatex
\documentclass[14pt,a4paper]{extarticle}

\usepackage{fontspec}
\setmainfont{Times New Roman}
\setsansfont{Arial}
\setmonofont{Consolas}
\usepackage[russian]{babel}
\usepackage{geometry}
\geometry{left=30mm,right=15mm,top=20mm,bottom=20mm}
\usepackage{setspace}
\onehalfspacing
\usepackage{amsmath,amssymb,mathtools}
\usepackage{booktabs,longtable,array,tabularx}
\usepackage{enumitem}
\usepackage{caption}
\usepackage{titlesec}
\usepackage{float}
\usepackage{xcolor}
\usepackage{fancyhdr}
\usepackage{lastpage}
\usepackage{microtype}
\usepackage{graphicx}

\usepackage[hidelinks]{hyperref}
\graphicspath{{figures/article_sources/}}

\setlength{\parindent}{12.5mm}
\setlength{\parskip}{0pt}
\setlength{\emergencystretch}{3em}
\setlist{nosep,leftmargin=12.5mm}
\renewcommand{\arraystretch}{1.18}
\renewcommand{\tablename}{Таблица}
\renewcommand{\figurename}{Рисунок}
\addto\captionsrussian{%
  \renewcommand{\refname}{СПИСОК ИСПОЛЬЗОВАННЫХ ИСТОЧНИКОВ}%
}
\titleformat{\section}[hang]
  {\normalfont\Large\bfseries\raggedright\hyphenpenalty=10000}
  {\thesection}{1em}{}
\titleformat{\subsection}[hang]
  {\normalfont\large\bfseries}{\thesubsection}{1em}{}
\titleformat{\subsubsection}[hang]
  {\normalfont\normalsize\bfseries}{\thesubsubsection}{1em}{}
\newcommand{\sectionbreak}{\clearpage}
\DeclareCaptionLabelSeparator{gostdash}{~---~}
\captionsetup{font=small,labelsep=gostdash}
\captionsetup[table]{position=top,skip=6pt,justification=raggedleft,
singlelinecheck=false}
\captionsetup[figure]{position=bottom,skip=6pt,justification=centering,
singlelinecheck=false}

\pagestyle{fancy}
\fancyhf{}
\fancyfoot[C]{\thepage}
\renewcommand{\headrulewidth}{0pt}

\newcommand{\source}[1]{\par\noindent\textit{Источник:} #1\par}
\newcommand{\phys}{\par\noindent\textbf{Физическая сущность.} }
\newcommand{\method}{\par\noindent\textbf{Методика.} }
\newcommand{\recommend}{\par\noindent\textbf{Рекомендация по выбору.} }
\newcommand{\subst}{\par\noindent\textbf{Подстановка.} }
\newcommand{\result}{\par\noindent\textbf{Результат.} }
\newcommand{\criterion}{\par\noindent\textbf{Критерий решения.} }
\newcommand{\stepgoal}[1]{\par\noindent\textbf{Цель шага.} #1\par}
\newcommand{\dd}{\mathop{}\!\mathrm{d}}
\newenvironment{whereblock}
 {\par\noindent где\par
  \begin{list}{}{\setlength{\leftmargin}{12.5mm}
                  \setlength{\itemindent}{-12.5mm}
                  \setlength{\itemsep}{0pt}
                  \setlength{\parsep}{0pt}
                  \setlength{\topsep}{0pt}}}
 {\end{list}}
\newcommand{\whereitem}[2]{\item[] #1 --- #2}


\usepackage{tikz}
\usepackage{xurl}
\setlength{\tabcolsep}{5pt}
\renewcommand{\theequation}{\arabic{equation}}




\title{Методика расчёта, проектирования и изготовления\\планарных силовых дросселей и трансформаторов\\с печатными обмотками}
\author{Самостоятельная инженерная методика\\Без привязки к конкретному изделию}
\date{2026}
\begin{document}
\maketitle\thispagestyle{empty}\newpage
\section*{РЕФЕРАТ}
\addcontentsline{toc}{section}{РЕФЕРАТ}
Настоящая методика устанавливает полный расчётный и производственный маршрут планарного силового магнитного компонента: электрический режим, выбор материала и витков, магнитная цепь и зазор, печатный рисунок и стек, потери, паразитные параметры, изоляция, нагрев, изготовление и испытания. Формулы и определения изложены в самом документе. Потокосцепление обозначается $\psi$.

Расчётные модели сопровождаются ограничениями. Документ не является протоколом испытанного изделия или подтверждением соответствия производства стандарту. Особое внимание уделено переходу от числа витков и требуемого сечения к реальной печатной дорожке.

\noindent\textbf{Ключевые слова:} планарный трансформатор, печатный дроссель, PCB, медная дорожка, стек слоёв, металлизированное отверстие, чередование обмоток, токораспределение, краевое поле зазора, изоляция, технологический допуск.
\newpage\tableofcontents\newpage
\section*{ТЕРМИНЫ, ОПРЕДЕЛЕНИЯ И СОКРАЩЕНИЯ}
\addcontentsline{toc}{section}{ТЕРМИНЫ, ОПРЕДЕЛЕНИЯ И СОКРАЩЕНИЯ}
\textbf{Планарная обмотка} --- обмотка из плоских проводников; в настоящем документе основным исполнением является медный рисунок одно- или многослойной печатной платы. Штампованные медные пластины и гибко-жёсткие платы рассматриваются как альтернативы, а не как автоматически эквивалентные технологии.

\textbf{Виток} --- проводящий путь, охватывающий рабочий магнитный поток. Число слоёв платы, число видимых колец и число витков электрической обмотки в общем случае различны.

\textbf{Секция} --- последовательно включённая часть обмотки. \textbf{Параллельная ветвь} --- полный путь между общими выводами, имеющий то же число витков и согласованную полярность, что и остальные ветви.

\textbf{Стек слоёв} (stack-up) --- последовательность медных и диэлектрических слоёв с указанием их готовых толщин, материалов и допусков. \textbf{Core платы} --- предварительно отверждённый ламинат; не путать с магнитным сердечником. \textbf{Prepreg} --- пропитанная связующим армирующая заготовка, используемая при прессовании платы.

\textbf{PTH} --- металлизированное сквозное отверстие. \textbf{NPTH} --- неметаллизированное отверстие. \textbf{DRC} --- проверка правил проектирования; \textbf{DFM} --- проверка технологичности. \textbf{Чередование обмоток} (interleaving) --- взаимное размещение секций разных обмоток по толщине платы.

\textbf{Функциональная, основная, дополнительная и усиленная изоляция} --- разные категории изоляции; требуемая категория устанавливается назначением изделия и применимым стандартом, а не толщиной паяльной маски. \textbf{Путь утечки} проходит по поверхности изоляции; \textbf{воздушный зазор изоляции} проходит по воздуху и не является магнитным зазором.

\section{НАЗНАЧЕНИЕ И ГРАНИЦЫ ПРИМЕНИМОСТИ}
\label{pl:scope}
\subsection{Назначение и физические функции}
\phys Планарность является способом изготовления обмотки, но не определяет электрическую функцию компонента. Дроссель накапливает энергию; трансформатор мостового или прямоходового преобразователя передаёт её между обмотками; обратноходовой магнитный компонент является накопительным связанным дросселем. Разделение этих функций принципиально при выборе зазора и тока, определяющего индукцию.\cite{web1,web35}

\method Для дросселя восстанавливают ток, потокосцепление, запасаемую энергию и дифференциальную индуктивность. Для трансформатора дополнительно определяют трансформацию, намагничивание, рассеяние, фазировку и изоляционный барьер. Нельзя подставить полный первичный ток нагрузки в однообмоточную формулу и объявить результат индукцией общего потока.

Частотная область применимости определяется данными материала, электрической длиной проводников и адекватностью сосредоточенной модели. При выраженных резонансах и существенных токах смещения необходима распределённая постановка. Планарная технология сама по себе не гарантирует малые потери или хорошую ЭМС.

\subsection{Разделение расчётных задач}
Дроссель выбирают по току, энергии и допустимой пульсации; передающий трансформатор --- по вольт-секундам, коэффициенту передачи и токам всех обмоток; обратноходовой компонент --- по накоплению и отдаче энергии с учётом разных интервалов проводимости. Во всех случаях после электромагнитного расчёта проверяют реализуемый медный рисунок, переходы, изоляцию и нагрев.

\criterion Принятое число витков и номинальная толщина меди ещё не являются проектом платы. Расчёт завершают только после получения согласованной геометрии и программы контроля.

\subsection{Правила использования источников}
Руководства TI используются как общая инженерная основа трансформаторов и высокочастотных обмоток; требования Eurocircuits --- как проверяемый пример реальных ограничений производства, но не как обязательные нормы любого завода. Полные тексты статей по PCB-дросселям и фольговым обмоткам использованы для проверки применимости полевых моделей.\cite{web1,web2,web17,web22,schaefer,ewaldfoil}

В расширенном списке литературы явно различаются полученный полный текст, доступная веб-страница и только проверенная библиографическая запись. Наличие DOI не означает, что статья прочитана целиком. Формулы из недоступных полных текстов не заимствованы. Новые геометрические и схемные соотношения ниже сопровождаются выводом или физическим основанием; они не выдаются за цитаты из таких статей.

\section{ИСХОДНЫЕ ДАННЫЕ И РАСЧЁТНАЯ ПОСЛЕДОВАТЕЛЬНОСТЬ}
\label{pl:inputs}
\subsection{Обязательные исходные данные}
До начала расчёта задают:
\begin{enumerate}
\item тип компонента и преобразователя; диапазоны входного и выходного напряжений, частоты, нагрузки и температуры; пуск, холостой ход, перегрузку и останов;
\item осциллограммы напряжения и тока каждой обмотки с направлениями отсчёта; для резонансного преобразователя --- рабочий диапазон частоты, циркулирующие токи и требуемые $L_m$, $L_\sigma$;
\item число изолированных цепей, рабочие и повторяющиеся пиковые напряжения между ними, импульсные воздействия, требуемую категорию изоляции;
\item габариты сердечника и окна с допусками, высоту PCB-пакета, форму отверстия под центральный стержень, расположение зазора, способ крепления и охлаждения;
\item согласованные с изготовителем готовые толщины меди и диэлектрика, ширину и зазор рисунка, допуск совмещения слоёв, размеры сверления, гарантированную металлизацию отверстий, допустимые размеры площадок и фрезерованных пазов;
\item свойства ламината: электрическая прочность по применимым условиям испытания, CTI при необходимости, температурные пределы, теплопроводность, $\varepsilon_r(f,T)$, $\tan\delta(f,T)$, требования к влаге и надёжности;
\item доступность контроля: микрошлифы, электрическая проверка голой платы, измерение сопротивлений, индуктивностей, ёмкостей, изоляции и температур внутренних слоёв.
\end{enumerate}
\recommend До выбора конкретного изготовителя платы не фиксируют универсальные значения минимального зазора, тока одного перехода и минимальной остаточной толщины prepreg. Такие значения являются входными данными проекта, а не константами методики.

\subsection{Расчётная последовательность и контрольные точки}
\begin{enumerate}
\item Классифицировать компонент и восстановить электрический режим.
\item Сформировать несколько вариантов сердечника, числа витков и технологии меди.
\item Для каждого варианта создать таблицу слоёв, соединений и потенциалов.
\item Разместить дорожки в CAD и проверить геометрию при допусках.
\item Рассчитать сопротивления и первичные оценки потерь, затем токораспределение, рассеяние и ёмкости.
\item Провести электромагнитно-тепловую итерацию и проверку изоляционного барьера.
\item Согласовать DFM и выпустить файлы изготовления.
\item Изготовить опытные платы, собрать магнитопровод, выполнить измерения.
\item Сопоставить измерение с моделью, уточнить модель и только затем утвердить конструкцию.
\end{enumerate}
\criterion Контрольная точка «расчёт» не заменяет точку «плата изготовлена», а испытанный единичный образец не заменяет квалификацию процесса. При отсутствии исходных данных результатом является параметризованная конструкция и перечень незакрытых ограничений, а не готовый к производству компонент.

\section{ОПРЕДЕЛЕНИЕ ЭЛЕКТРИЧЕСКОГО РЕЖИМА}
\label{pl:mode}
\subsection{Планарный дроссель и обратноходовой компонент}

\subsection{Напряжение, ток и потокосцепление}
\phys Потокосцепление обозначается строчной буквой $\psi$. Магнитный поток $\Phi$ и потокосцепление различают: для $N$ одинаково сцепленных витков $\psi=N\Phi$. Единица потокосцепления --- Вб (допускается пояснение Вб$\cdot$виток); число витков безразмерно. В многообмоточной модели $\psi_p$ и $\psi_s$ относятся к соответствующим выводам, а $\psi_m$ --- только к общему потоку.

\method Напряжение на индуктивной части одиночного дросселя и изменение потокосцепления определяют при пассивном направлении отсчёта:
\begin{equation}
 u_L=u_{\rm term}-u_R,\qquad u_R\simeq R_{\rm DC}(\vartheta)i.
 \label{eq:inductive-voltage}
\end{equation}
Здесь $u_{\rm term}$ --- клеммное напряжение, $u_R$ --- падение в проводнике, В; $i$ --- ток, А. Последняя оценка не учитывает частотно-зависимое внутреннее падение; при значительных AC-потерях используют соответствующую модель импеданса.
\begin{equation}
 u_L=\frac{d\psi}{dt},\qquad
 \Delta\psi_{12}=\int_{t_1}^{t_2}u_L(t)\,dt.
 \label{eq:flux-linkage-voltage}
\end{equation}
\begin{equation}
 \psi(t)=\psi(t_0)+\int_{t_0}^{t}u_L(\tau)\,d\tau.
 \label{eq:flux-linkage-trajectory}
\end{equation}
Начальное $\psi(t_0)$ определяют по магнитному состоянию, постоянному току и предыстории. Интегрирование осциллограммы без этой константы не доказывает отсутствие насыщения. Для общего потока трансформатора вместо полного $u_L$ интегрируют ЭДС общей ветви, исключив составляющую рассеяния.
\begin{equation}
 L_{\rm diff}=\left.\frac{\partial\psi}{\partial i}\right|_{\vartheta},\qquad
 \frac{di}{dt}=\frac{u_L}{L_{\rm diff}(i,\vartheta)}.
 \label{eq:current-nonlinear}
\end{equation}
Это квазистатическая однозначная модель при практически постоянной за период температуре; при гистерезисе и динамике потерь требуется модель состояния. Для линейного $L$ изменение тока равно $\Delta\psi/L$.
\begin{equation}
 I_{\rm avg}=\frac1{T_s}\int_0^{T_s}i(t)\,dt,\quad
 I_{\rm rms}=\sqrt{\frac1{T_s}\int_0^{T_s}i^2(t)\,dt},\quad I_{\rm pk}=\max_t|i(t)|.
 \label{eq:current-general}
\end{equation}
При треугольной пульсации вокруг неотрицательного среднего тока:
\begin{equation}
 I_{\max}=I_{\rm avg}+\frac{\Delta I_{pp}}2,\qquad
 I_{\rm rms}=\sqrt{I_{\rm avg}^2+\frac{\Delta I_{pp}^2}{12}}.
 \label{eq:ipeak}
\end{equation}
Здесь $T_s$ --- период, с; $\Delta I_{pp}$ --- размах тока, А. При паузах нулевого тока, резонансном токе или обратном токе применяют \eqref{eq:current-general}, а не треугольную формулу автоматически.
\begin{equation}
 W=\frac12 LI^2.
 \label{eq:energy-linear}
\end{equation}
\begin{equation}
 W(i,\vartheta)=\int_0^{\psi(i,\vartheta)}i(\widetilde\psi,\vartheta)\,d\widetilde\psi.
 \label{eq:energy-nonlinear}
\end{equation}
$W$ --- энергия, Дж. Второе выражение предполагает однозначную безгистерезисную характеристику и согласованное нулевое состояние. Интеграл $\int\psi\,di$ является коэнергией, а не той же энергией при нелинейности. Энергия дросселя не заменяется передаваемой мощностью трансформатора.\cite{web1,web35}
\begin{equation}
 B_e(t)=\frac{\psi_m(t)}{NA_e},\qquad
 \Delta B_{pp}=B_{e,\max}-B_{e,\min}.
 \label{eq:b-equivalent-trajectory}
\end{equation}
$A_e$ --- эффективное сечение, м$^2$; $B_e$ --- средняя индукция, Тл. Не подставляют полное потокосцепление выводов с заметным рассеянием вместо $\psi_m$. В неоднородном сердечнике проверяют также минимальное сечение и локальные поля.
\begin{equation}
 \max_{r,t,\mathbf x}|B_r(\mathbf x,t)|\leq B_{\rm allow}(\vartheta_{\max}).
 \label{eq:b-acceptance}
\end{equation}
Индекс $r$ обозначает режим и сочетание допусков. $B_{\rm allow}$ задают по материалу, температуре, допустимой нелинейности и запасу, а не универсальной долей насыщения. Потери дополнительно могут потребовать меньший размах.\cite{web3}


Обратноходовой компонент рассчитывают как связанный накопительный дроссель. На интервалах накопления и отдачи энергии активны разные обмотки; ток меди и ток намагничивания имеют разные формы. Число витков вторичной обмотки выбирают вместе с отражённым напряжением, напряжениями ключа и выпрямителя и схемой ограничения выбросов. Формулы мостового трансформатора не переносят на flyback без изменения модели.\cite{web35}

\subsection{Общий поток и коэффициент трансформации}
\phys Для идеальной связи ЭДС витка одинакова в каждой обмотке, охватывающей один и тот же поток. Поэтому при согласованной фазировке
\begin{equation}
 a=\frac{N_p}{N_s}=\frac{e_p}{e_s},\qquad
 N_p i_p+N_s i_s=\mathcal F_m .
 \label{pl:eq:ratio-mmf}
\end{equation}
\begin{whereblock}
\whereitem{$N_p,N_s$}{числа витков первичной и вторичной обмоток;}
\whereitem{$e_p,e_s$}{ЭДС общего потока, В, не обязательно клеммные напряжения;}
\whereitem{$i_p,i_s$}{токи, положительные при входе в одноимённые начала обмоток, А;}
\whereitem{$\mathcal F_m$}{результирующая намагничивающая сила, А.}
\end{whereblock}
Соотношение получено из закона индукции и закона полного тока. При передаче мощности ток вторичной обмотки с выбранным знаком компенсирует большую часть первичных ампер-витков. Намагничивающий ток не равен полному первичному току.\cite{web1}

\method Для $e_p(t)$ используют связь напряжения с потокосцеплением из \eqref{eq:flux-linkage-voltage} и траекторию из \eqref{eq:flux-linkage-trajectory}. Измеренное клеммное напряжение очищают от падения в меди и рассеяния. Для выбранного $N_p$ восстанавливают $B(t)$ согласно \eqref{eq:b-equivalent-trajectory}, находят размах по \eqref{eq:b-equivalent-trajectory} и проверяют \eqref{eq:b-acceptance}. Интегрирование даёт изменение потока, но не определяет неизвестную начальную индукцию; пуск, остаточное намагничивание и смещение потока рассматривают отдельно.

Коэффициент $a$ не выбирают только как отношение постоянных входного и выходного напряжений: учитывают топологию, коэффициент заполнения, падения напряжения, коэффициент передачи резонансного контура и пределы регулирования. У прямоходового преобразователя проверяют отдельный интервал размагничивания, у мостового --- симметрию вольт-секунд и защиту от ухода потока.

\subsection{Намагничивание, рассеяние и эквивалентная схема}
В линейной малосигнальной модели двух обмоток с согласованными началами:
\begin{equation}
 \begin{bmatrix}\psi_p\\\psi_s\end{bmatrix}
 =\begin{bmatrix}L_{pp}&M\\M&L_{ss}\end{bmatrix}
 \begin{bmatrix}i_p\\i_s\end{bmatrix},\qquad
 k=\frac{M}{\sqrt{L_{pp}L_{ss}}}.
 \label{pl:eq:matrix-L}
\end{equation}
Здесь $\psi$ --- потокосцепление, Вб; $L_{pp},L_{ss}$ --- собственные индуктивности при разомкнутой другой обмотке, Гн; $M$ --- взаимная индуктивность, Гн. Для пассивной взаимной линейной модели $|k|\leq1$. При насыщении вместо постоянной матрицы необходима зависимость потокосцеплений от всех токов либо её дифференциальная матрица около рабочей точки.

При идеальном коротком замыкании вторичной обмотки, пренебрежимо малых сопротивлениях и ёмкостях из условия нулевого приращения вторичного потокосцепления следует
\begin{equation}
 L_{\mathrm{sc},p}=L_{pp}-\frac{M^2}{L_{ss}}
                 =L_{pp}(1-k^2).
 \label{pl:eq:short-circuit}
\end{equation}
Это входная индуктивность короткого замыкания, а не универсальная идентификация отдельно первичного рассеяния. При конечном сопротивлении перемычки и вторичной меди извлекают полную комплексную модель; измеренное $\operatorname{Im}Z/\omega$ может существенно отличаться от идеального выражения.

\recommend Для LLC и других резонансных преобразователей указывают, какая именно индуктивность является параметром схемы и к какой стороне она приведена. «Минимум рассеяния» не является универсальной целью: рассеяние может быть частью заданного резонансного контура.\cite{web1}

\section{ВЫБОР МАРКИ МАТЕРИАЛА, ТИПОРАЗМЕРА, ЧИСЛА ВИТКОВ И ЗАЗОРА}
\label{pl:core}

\subsection{Материал, число витков и магнитная цепь}
\method Выбирают материал по данным изготовителя: $B(H,\vartheta)$, комплексной проницаемости, потерям при требуемой частоте и форме индукции, температурному диапазону, допускам и доступности. Геометрические $A_e,l_e,V_e,A_{\min}$ и размеры окна получают для конкретного типоразмера. Каталожные данные без указания температуры и амплитуды недостаточны.\cite{web3}

Для простого неразветвлённого потока линейная магнитная цепь даёт
\begin{equation}
 \mathcal R_m=\sum_j\frac{l_j}{\mu_0\mu_{r,j}A_j}
 +\sum_q\frac{g_q}{\mu_0F_q A_{g,q}},\qquad
 L_m=\frac{N^2}{\mathcal R_m},\qquad A_L=\frac{L_m}{N^2}.
 \label{eq:magnetic-circuit}
\end{equation}
$\mathcal R_m$ --- магнитное сопротивление, А/Вб; $l_j,A_j,\mu_{r,j}$ --- длина, сечение и относительная проницаемость участка сердечника; $g_q,A_{g,q}$ --- длина и сечение зазора; $F_q$ --- определённый для этой модели коэффициент увеличения эффективного сечения зазора. Для разветвлённого сердечника сначала решают сеть потоков: простая сумма всех внешних зазоров может быть неверна. Нелинейную задачу решают по $H(B)$, а не подстановкой начальной проницаемости во всём диапазоне.

Для одного зазора при известном беззазорном $A_{L0}$:
\begin{equation}
 g^{(n+1)}=\mu_0 A_g F(g^{(n)})
 \left(\frac1{A_{L,\rm target}}-\frac1{A_{L0}}\right).
 \label{eq:gap-fringing}
\end{equation}
Для нескольких последовательных зазоров в том же потоке:
\begin{equation}
 \sum_q\frac{g_q}{F_q A_{g,q}}=
 \mu_0\left(\frac1{A_{L,\rm target}}-\frac1{A_{L0}}\right).
 \label{eq:multiple-gap-balance}
\end{equation}
Это предварительные оценки. Коэффициенты $F_q$ зависят от геометрии; в нулевом приближении можно принять $F_q=1$ с явным указанием пренебрежения выпучиванием. Если требуемый $A_L$ выше беззазорного, положительный зазор решение не даст.

Для дросселя перебирают целое $N$, вычисляют требуемый $A_L=L/N^2$, подбирают зазор и проверяют $L_{\rm diff}(i,\vartheta)$, насыщение, потери, высоту и ширину обмотки. Для передающего трансформатора сначала выбирают $N_p,N_s$ по коэффициенту передачи и вольт-секундам, затем проверяют намагничивание и размещение. По размаху ЭДС общей ветви требуется
\begin{equation}
 N_p\geq\frac{\max_t S(t)-\min_t S(t)}{A_e\Delta B_{\rm allow}},
 \qquad S(t)=\int_{t_0}^t e_p(\tau)\,d\tau.
 \label{eq:turns-volts}
\end{equation}
Целое число округляют вверх, после чего отдельно проверяют абсолютный уровень индукции и пуск. Один периодический размах не ограничивает смещение потока.

\subsection{Потери сердечника и область применимости модели}
Для синусоидальной индукции допускается аппроксимация $p_v=k f^\alpha\widehat B^\beta$ в согласованных единицах; $\widehat B=\Delta B_{pp}/2$. Для одного главного цикла без вложенных петель используют форму iGSE только с верифицированными коэффициентами:
\begin{equation}
 p_v=\frac1{T_s}\int_0^{T_s}k_i\left|\frac{dB}{dt}\right|^\alpha
 (\Delta B_{pp})^{\beta-\alpha}\,dt.
 \label{eq:igse}
\end{equation}
Здесь $p_v$ --- Вт/м$^3$, $k_i,\alpha,\beta$ относятся к данному материалу, температуре и частотному диапазону. $k_i$ не равен автоматически синусоидальному $k$. Согласование с синусоидой даёт
\begin{equation}
 k_i=\frac{k}{2^{\beta-\alpha}(2\pi)^{\alpha-1}
 \int_0^{2\pi}|\cos\theta|^\alpha d\theta}.
 \label{eq:igse-coefficient}
\end{equation}
Это алгебраическая нормировка выбранной модели, не доказательство её точности для любого феррита. При DC-подмагничивании, вложенных петлях и релаксации нельзя объявлять поправку известной: требуются измеренные зависимости либо другая подтверждённая модель. При отсутствии данных расчёт потерь считается незакрытым.
\begin{equation}
 P_{\rm Fe}=\int_{V_{\rm Fe}}p_v(\mathbf x)\,dV\simeq p_v V_e.
 \label{eq:core-loss-volume}
\end{equation}
Последнее равенство --- однородная каталожная оценка, а не локальный полевой результат. Нельзя переносить коэффициенты другого материала или смешивать мТл и Тл, кГц и Гц.


\subsection{Планарные габаритные ограничения}
Для пакета платы в окне сердечника вводят геометрическое условие
\begin{equation}
 h_{\mathrm{stack,max}}+h_{\mathrm{ins,max}}+h_{\mathrm{assembly,max}}
 \leq h_{\mathrm{window,min}}.
 \label{pl:eq:height}
\end{equation}
Здесь $h_{\mathrm{stack,max}}$ --- максимальная готовая толщина PCB-пакета; $h_{\mathrm{ins,max}}$ --- суммарная толщина дополнительных изоляционных прокладок, не включённых в PCB; $h_{\mathrm{assembly,max}}$ --- монтажный запас, клей и неплоскостность; $h_{\mathrm{window,min}}$ --- минимальная высота окна с учётом допусков. Все величины задают в одной системе единиц. Маску учитывают один раз: либо в толщине готового пакета, либо отдельным слагаемым.

Отдельно проверяют отверстие под центральный стержень, допуск фрезерования, радиусы внутренних углов, зазор к хрупкому ферриту и электрическую изоляцию от него. Плата, помещающаяся в окно по номиналу, может не собраться при неблагоприятном сочетании допусков. Нельзя изгибать плату или стягивать сердечник до её деформации, чтобы компенсировать ошибку высоты.

\subsection{Воздушный зазор и расстояние до меди}
\phys Краевое поле зазора проходит через широкую медь и вызывает вихревые токи; сама медь экранирует поле и может менять индуктивность. Одномерная магнитная цепь этого не описывает.\cite{ewaldfoil,schaefer}

\method Применяют \eqref{eq:gap-fringing} и \eqref{eq:multiple-gap-balance} как предварительные оценки. Затем сравнивают центральный зазор, распределение нескольких зазоров, отступ меди и альтернативное расположение зазора. Минимальный отступ выбирают по совместной зависимости потерь, температуры и индуктивности от геометрии, а не по универсальному правилу вида «несколько толщин меди».

Статья о компенсации краевого поля показывает возможность улучшения токораспределения специальной геометрией; это не разрешение произвольно приближать медь к любому зазору.\cite{schaefer} Каждое такое решение проверяют на реальном стеке и возбуждении.

\criterion Коэффициент выпучивания, идентифицированный из FEM, относится только к геометрии и возбуждению данной модели. Если проводящая плата заметно экранирует поле, частотно-зависимую индуктивность нельзя описывать одной постоянной поправкой к статическому зазору.

\recommend В работе Schäfer--Bortis--Kolar рассмотрена компенсация краевого поля специальным положением зазора. Для прямой тонкой полосы и для кругового витка оптимальные положения различаются; при круговом витке даже DC-плотность тока неравномерна. Перенос на многослойный пакет основан на малом расстоянии между слоями относительно ширины полосы. При нарушении этого условия требуется полная полевая задача. Авторы отдельно выделяют потери выводов и рекомендуют копланарное подключение; сквозные переходы выбраны по стоимости, а не как физический запрет blind/buried via.\cite[разделы II--V]{schaefer}

\section{ПРОЕКТИРОВАНИЕ ПЕЧАТНОЙ ПЛАТЫ И ДОРОЖКИ ОБМОТКИ}
\label{pl:pcb}
\subsection{Выбор способа изготовления обмотки}
\begin{longtable}{>{\raggedright\arraybackslash}p{34mm}>{\raggedright\arraybackslash}p{60mm}>{\raggedright\arraybackslash}p{60mm}}
\caption{Сравнение технологий}\label{pl:technologies}\\
\toprule Вариант & Преимущества & Ограничения\\\midrule\endfirsthead
\toprule Вариант & Преимущества & Ограничения\\\midrule\endhead
Обычная многослойная PCB & Повторяемый рисунок, интеграция с силовой схемой, автоматизированный контроль & Тонкая медь, потери в переходах, значительные межслойные ёмкости\\
PCB с толстой медью & Большее DC-сечение при той же ширине & Более грубые зазоры и рисунок, заполнение смолой, усложнение прессования; толщина может усилить AC-потери\\
Пакет отдельных плат & Возможность менять секционирование и изоляцию & Дополнительные соединения, допуски пакета, контактные потери\\
Медные пластины / фольга с PCB & Большое сечение сильноточных витков & Отдельные операции изоляции, фиксации и контроля кромок; не обычный PCB-процесс\\
Печатный литцендрат & Возможность транспозиции ветвей & Много слоёв и переходов, сниженное заполнение медью; эффект нужно доказать\\
\bottomrule
\end{longtable}
\recommend Сначала сравнивают простые технологические варианты. Печатный литцендрат и HDI не вводят без подтверждённого выигрыша в потерях, габарите или надёжности. Набор параллельных дорожек без транспозиции не становится литцендратом.

\subsection{Шаг 1. Создать электрическую карту витков до трассировки}
\stepgoal{Получить однозначную связь между электрической схемой и медными слоями.}
В таблице каждого слоя указывают имя слоя CAD, обмотку, секцию, последовательный номер витка, направление обхода сердечника, начало и конец, потенциал относительно изоляционного барьера, толщину меди и переход к следующему слою. Для параллельных ветвей отдельно задают одинаковое число последовательных витков.

Например, два слоя по два последовательных витка, соединённые согласно, образуют четырёхвитковую обмотку. Две такие четырёхвитковые ветви в параллели по-прежнему образуют четыре, а не восемь витков. При встречном последовательном соединении вклад слоёв вычитается. При встречном параллельном соединении возникает недопустимый контур с большой циркулирующей составляющей тока.

\method Прослеживают непрерывный путь от входного до выходного вывода и направление обхода центрального стержня. Нижний слой оценивают в общей системе координат платы, а не по визуальному впечатлению от зеркального вида. Фазировку проверяют отдельным малосигнальным испытанием до включения силового преобразователя.

\subsection{Шаг 2. Построить механическую основу и зоны запрета меди}
\begin{enumerate}
\item Импортировать проверенные размеры сердечника и создать контур PCB и вырез под центральный стержень.
\item Учесть радиус фрезы и максимальный габарит стержня. Вырез задавать замкнутым контуром, а не цепочкой перекрывающихся сверлений.
\item Создать отдельные зоны запрета для меди у торца платы, магнитного зазора, первично-вторичного барьера, крепежа и выводов.
\item Зарезервировать коридор выводов и площадок переходов: его нельзя считать доступной шириной для равномерной спирали.
\item Проверить доступный рисунок на всех слоях и при допустимом смещении сверления и слоёв.
\end{enumerate}
Технологический отступ меди от кромки платы и безопасный отступ изоляции имеют разное назначение; принимают более строгое требование. Значения одного завода не распространяются автоматически на другую технологию.\cite{web17,web18,web21}


\subsection{Расчётные выражения для меди и переходов}
\begin{equation}
 A_{\rm req}=\frac{I_{\rm rms}}{J_0},\qquad
 w_0=\frac{A_{\rm req}}{n_\parallel t_{\min}}.
 \label{eq:copper-area}
\end{equation}
$J_0$ --- предварительная проектная плотность, А/м$^2$; $n_\parallel$ --- число равномерно нагруженных полных параллельных ветвей; $t_{\min}$ --- гарантированная готовая медь, м. $J_0$ не является универсальной допустимой нормой: кандидат принимается только после теплового расчёта.
\begin{equation}
 \rho(\vartheta)=\rho_{20}[1+\alpha_{\rm Cu}(\vartheta-20\,^{\circ}\mathrm C)].
 \label{eq:rho-temperature}
\end{equation}
Для иллюстративного расчёта приняты $\rho_{20}=1{,}724\cdot10^{-8}$ Ом$\cdot$м и $\alpha_{\rm Cu}=3{,}93\cdot10^{-3}$ К$^{-1}$. Это заданные номинальные параметры стандартной меди, не гарантированные характеристики конкретной фольги или гальваники; для выпуска их заменяют паспортными или измеренными значениями.
\begin{equation}
 R_b=\sum_k\rho_k\frac{l_k}{A_k}+R_{\rm via,b}+R_{\rm joints,b},
 \qquad R_{\rm eq}=\left(\sum_b\frac1{R_b}\right)^{-1}.
 \label{eq:planar-rdc}
\end{equation}
$b$ --- полная параллельная ветвь, $k$ --- её последовательный участок; $A_k$ --- фактическое сечение, $l_k$ --- длина; добавлены переходы и соединения. Формула параллели справедлива при общих выводах и DC, без источников наведённой ЭДС. Сужения считают отдельно; при неоднородном DC-поле вместо $l/A$ решают задачу проводимости.
\begin{equation}
 A_v=\frac\pi4[(d_f+2t_v)^2-d_f^2],\quad
 R_v=\frac{\rho l_v}{A_v},\quad R_{\rm group}\simeq\frac{R_v}{n_v}.
 \label{eq:via-resistance}
\end{equation}
$d_f$ --- готовый внутренний диаметр; $t_v$ --- минимальная медь стенки; $l_v$ --- участок между подключёнными слоями, м. Последнее выражение предполагает одинаковые переходы и равномерную нагрузку, не учитывает сопротивление растекания в площадках.

\subsection{Шаг 3. Выбрать ширину по минимальному готовому сечению}
Первичную оценку требуемого сечения получают по \eqref{eq:copper-area}, а ширины --- по \eqref{eq:copper-area}. Это только начальный кандидат. Для PCB в расчёт подставляют \emph{готовую минимальную} толщину, а не номинальную массу фольги в унциях.

Если изготовитель задаёт абсолютное возможное уменьшение ширины $\Delta w_-$, то для выбранной CAD-ширины $w_{\mathrm{CAD}}$ минимальная ширина равна
\begin{equation}
 w_{\min}=w_{\mathrm{CAD}}-\Delta w_-,\qquad
 A_{\mathrm{Cu,min}}=t_{\min}w_{\min}.
 \label{pl:eq:minimum-section}
\end{equation}
Это прямоугольная консервативная модель сечения. Если ширина после травления существенно различается сверху и снизу и обе величины подтверждены микрошлифом, вместо прямоугольника используют геометрию трапеции:
\begin{equation}
 A_{\mathrm{Cu}}=t_{\mathrm{Cu}}\frac{w_{\mathrm{top}}+w_{\mathrm{bottom}}}{2}.
 \label{pl:eq:trapezoid}
\end{equation}
Здесь ширины и толщина относятся к одному фактическому сечению. Нельзя совместить несогласованные средние толщину и ширину разных образцов и объявить их гарантированным минимумом.

\source Производственные причины различия начальной и конечной толщины, травления боковых стенок и разброса гальваники рассмотрены в \cite{web16,web22}. Соотношения \eqref{pl:eq:minimum-section} и \eqref{pl:eq:trapezoid} следуют из геометрии, а не являются нормативными допусками.

\subsection{Шаг 4. Проверить размещение дорожек в окне}
Для $n$ параллельно расположенных участков витков одинаковой ширины в наиболее узком сечении доступной полосы:
\begin{equation}
 n w_{\mathrm{CAD}}+(n-1)s_{\mathrm{CAD}}+m_{\mathrm{in}}+m_{\mathrm{out}}
 \leq b_{\mathrm{available}}.
 \label{pl:eq:radial-fit}
\end{equation}
\begin{whereblock}
\whereitem{$n$}{число соседних проводящих полос в данном сечении, не число параллельных электрических ветвей;}
\whereitem{$s_{\mathrm{CAD}}$}{проектный междорожечный зазор;}
\whereitem{$m_{\mathrm{in}},m_{\mathrm{out}}$}{резервы по внутренней и наружной границам;}
\whereitem{$b_{\mathrm{available}}$}{доступная ширина от механической геометрии в том же сечении.}
\end{whereblock}
Условие проверяют на обеих сторонах окна, в углах, у переходов и выводов. Для трапецеидального или криволинейного окна расчёт одного среднего сечения недостаточен.

Междорожечный зазор выбирают по более строгому из технологического и электрического требований. Для электрической проверки находят максимальное напряжение между фактически соседними участками: у выводов рядом могут оказаться начало и конец всей обмотки. При быстрых фронтах распределение напряжения может быть неравномерным, поэтому деление полного напряжения на число витков без проверки не задаёт безопасный зазор.

Если требуемая по потерям ширина не помещается, не уменьшают зазор произвольно. Последовательно сравнивают меньшее число витков на слой и большее число последовательных слоёв; дополнительные полные параллельные ветви; большее окно сердечника; другую технологию меди. Каждый вариант повторно проверяют по высоте, числу переходов, AC-потерям и изоляции.

Проверка по максимальному контуру меди и проверка по минимальному сечению решают разные задачи: максимальная ширина опасна для изоляционного зазора, минимальная --- для потерь и нагрева. Компенсацию травления согласуют с заводом; не прибавляют её к Gerber повторно, если изготовитель уже выполняет компенсацию в CAM.

\subsection{Шаг 5. Нарисовать непрерывную дорожку}
\stepgoal{Получить последовательную обмотку без замкнутых короткозамкнутых колец и случайных перемычек.}
\begin{enumerate}
\item Назначить оба вывода обмотки и межслойные соединения на электрической схеме. Для многовитковой обмотки вести таблицу внутренних узлов секций.
\item На выбранном слое провести непрерывный обход стержня дорожкой постоянной расчётной ширины либо контролируемой областью меди.
\item Между соседними обходами оставить рассчитанный изоляционный зазор. У выводов сохранить разрыв между началом и концом витка.
\item Использовать плавные или скошенные повороты; проверять минимальную ширину по нормали к пути тока. Внутренний угол не должен создавать узкую перемычку.
\item Вывести внутренний конец через предусмотренную группу переходов на следующий слой. Переход на том же слое поперёк витков недопустим.
\item Проверить путь тока по всем слоям; ширину, длину и узкие места измерять по фактической геометрии, а не по условному среднему диаметру.
\end{enumerate}

\begin{figure}[H]\centering
\begin{tikzpicture}[scale=0.75,>=stealth]
\fill[gray!20] (-1.6,-1.2) rectangle (1.6,1.2);
\draw (-1.6,-1.2) rectangle (1.6,1.2);
\node[align=center,font=\small] at (0,0) {Стержень\\и изоляционный отступ};
\draw[line width=6pt,brown!70!black] (3,-1.5)--(3,2.3)--(-3,2.3)--(-3,-2.3)--(3,-2.3);
\draw[->,line width=1pt] (-1.3,2.8)--(1.3,2.8);
\node[right,font=\small] at (3,-1.5) {конец B};
\node[right,font=\small] at (3,-2.3) {начало A};
\draw[<->] (2.4,-1.5)--(2.4,-2.3);
\node[font=\small,align=center] at (0,-3.3) {Разрыв между A и B обязателен.\\Замыкание разрыва создаёт короткозамкнутый виток.};
\end{tikzpicture}
\caption{Принцип рисунка одного печатного витка; схема не является масштабным Gerber-шаблоном}
\label{pl:fig:turn}
\end{figure}

\recommend Если обмотка оформлена одним net, стандартный DRC может разрешать слияние её витков как меди одной цепи. Поэтому включают проверку зазоров внутри одной цепи, если CAD её поддерживает, либо разделяют секции контролируемыми внутренними узлами/net-tie и проверяют итоговую связность отдельно. Отсутствие ошибок обычного DRC ещё не доказывает правильное число витков.

\subsection{Шаг 6. Проверить повороты, выводы и сужения}
\phys Внутренние углы, площадки переходов и выводы создают неравномерную плотность тока даже при постоянной толщине меди. Широкая дорожка не компенсирует узкое горлышко на входе.

\method Каждое сужение включают отдельным участком в \eqref{eq:planar-rdc}. Для переменного тока выделяют его в полевой модели. Повороты делают с возможным по габариту внутренним радиусом; универсальное отношение радиуса к ширине без проверки не назначают. Тепловые спицы у силовых площадок рассматривают как реальные параллельные проводники и включают в расчёт. Сплошное подключение лучше проводит ток, но требует проверки паяемости и профиля пайки.

Отдельно проверяют длину и ширину выводов вне окна: они могут определять существенную часть сопротивления и рассеяния. Для сильноточных площадок допускают усиление медной пластиной только как отдельный технологический вариант с контролем соединения, изоляции и термоциклирования. Нанесение припоя на дорожку не заменяет гарантированное сечение меди.

\subsection{Шаг 7. Выбрать переходы между слоями}
Сопротивление металлизированных цилиндров считают по \eqref{eq:via-resistance}. Минимальная толщина стенки, готовый диаметр отверстия и реальная длина проводящего участка должны быть заданы изготовителем. Для межслойного перехода ток проходит участок между подключёнными слоями; не всегда всю толщину платы. Остальная часть сквозного цилиндра может влиять на паразитные параметры.

Группа переходов должна принимать ток по всей доступной ширине дорожки, а не через узкую шейку. Размещение всех переходов в длинный ряд вдоль направления тока может нагружать ближайшие отверстия сильнее дальних. Параллельное деление сопротивления на число отверстий допустимо только как предварительная равномерная модель.

Для тепловой проверки вводят коэффициент неравномерности
\begin{equation}
 k_{\mathrm{via}}=\frac{\max_r I_{r,\mathrm{rms}}}{I_{\mathrm{group,rms}}/n_v}.
 \label{pl:eq:via-imbalance}
\end{equation}
Здесь $n_v$ --- число переходов; токи получены при фактическом подключении и возбуждении. Это определение, а не заданная поправка. Если знаменатель близок к нулю при циркулирующих токах, оценивают непосредственно каждый $I_{r,\mathrm{rms}}$, а не отношение.

\criterion Нельзя применять универсальное правило «один via выдерживает один ампер». Проверяют цилиндр, контактное кольцо, подвод меди, температурный предел и устойчивость к термоциклам. Внутренние отверстия могут быть существенно горячее доступной для тепловизора поверхности.

\source Заводские правила различают готовый размер и инструмент сверления; автоматическое уменьшение отверстий в CAM может изменить расчётные параметры. Для критических переходов запрещают несогласованные изменения и запрашивают утверждение производственных данных.\cite{web18}

\subsection{Шаг 8. Спроектировать стек слоёв}
\begin{longtable}{>{\raggedright\arraybackslash}p{18mm}>{\raggedright\arraybackslash}p{42mm}>{\raggedright\arraybackslash}p{94mm}}
\caption{Пример функциональной карты четырёх медных слоёв, не готовый стек}\label{pl:stack-example}\\
\toprule Слой & Назначение & Обязательная проверка\\\midrule\endfirsthead
\toprule Слой & Назначение & Обязательная проверка\\\midrule\endhead
L1 & Первичная секция P1 & Начало, конец, число витков, минимальная конечная медь\\
D12 & Изоляция P--S & Категория барьера, остаточная толщина, рабочее напряжение\\
L2 & Вторичная ветвь S1 & Полярность, ток, симметрия к S2\\
D23 & Межслойная изоляция & Напряжение между S1 и S2 не принимают нулевым без проверки переходных режимов\\
L3 & Вторичная ветвь S2 & Если параллельно S1: те же витки и полярность; не складывать их число витков\\
D34 & Изоляция S--P & Независимая проверка второго интерфейса\\
L4 & Первичная секция P2 & Последовательное согласное соединение с P1; длина переходов\\
\bottomrule
\end{longtable}
Структура P1--S1--S2--P2 является только кандидатом. Преимущество по рассеянию не отменяет увеличения площади первично-вторичного перекрытия. Для высоковольтного барьера такой стек может оказаться хуже разнесённого.

В файле изготовления указывают наименование материала, медь каждого слоя, тип core/prepreg, готовые толщины и допуски. Обозначения «FR-4» и «плата 1,6 мм» не задают полный стек и не определяют остаточную межслойную изоляцию. Учитывают заполнение смолой рядом с толстой медью и локальное уменьшение диэлектрика после прессования. Симметрию пакета и баланс меди согласуют с заводом; дополнительные балансирующие полигоны в магнитном поле нельзя добавлять автоматически.\cite{web19,web20,web22}

\subsection{Шаг 9. Соединить параллельные дорожки без неучтённых контуров}
\method Каждая параллельная ветвь должна иметь одинаковое число витков, одинаковую полярность и сопоставимые полные импедансы. Равенство DC-длин является необходимым полезным условием, но не гарантирует равного тока при разных положениях слоёв в магнитном поле.

Соединяют ветви в заранее выбранных узлах. Произвольное сшивание промежуточных точек может создавать контуры между слоями, охватывающие поток рассеяния. Транспозицию выполняют по явной карте переходов, где каждая ветвь последовательно занимает соответствующие позиции. Стоимость транспозиции --- дополнительная длина, переходы и места сужений; её выигрыш проверяют по полным потерям.

\criterion Для каждой ветви отчёт содержит ток, DC- и AC-потери, максимальную температуру и схему соединений. Простое умножение сечения одного слоя на число слоёв не является достаточной проверкой.

\subsection{Матричная проверка параллельных ветвей}
Для фиксированных геометрии, частоты и линейного состояния извлекают матрицу импедансов раздельных слоёв. Матрица соединения $\mathbf C_s$ задаёт, какие слои последовательно принадлежат каждой ветви, включая знак фазировки:
\begin{equation}
 \mathbf Z_b=\mathbf C_s^T\mathbf Z_{\ell}\mathbf C_s,\qquad
 \mathbf i_b=\mathbf Z_b^{-1}\mathbf 1\,U.
 \label{eq:parallel-branch-matrix}
\end{equation}
$\mathbf Z_{\ell}$ и $\mathbf Z_b$ --- комплексные матрицы импедансов слоёв и ветвей, Ом; $U$ --- общее напряжение; $\mathbf 1$ --- столбец единиц. Равенство комплексных сумм строк $\mathbf Z_b^{-1}$ означает равные токи при общем напряжении. При заданном общем токе:
\begin{equation}
 \mathbf i_b=\frac{\mathbf Z_b^{-1}\mathbf 1}{\mathbf 1^T\mathbf Z_b^{-1}\mathbf 1}I_{\rm total}.
 \label{eq:parallel-fixed-current}
\end{equation}
Матрица должна учитывать поле зазора. Этот подход позволяет сравнивать соединения после извлечения матрицы, не повторяя FEM для каждой чисто электрической перестановки. Если меняются выводы, переходы, геометрия или нелинейная рабочая точка, матрицу обновляют.\cite[разделы IV--V]{doi35}

Равенство токов и минимум полных потерь --- разные критерии. В рассмотренном авторами EI-сердечнике сплошное чередование слоёв не устраняет асимметрию относительно зазора. Сравнивать потери нужно при одинаковом заданном токе или при одинаковом напряжении, не смешивая эти постановки. В методике используется формула потерь с действующими комплексными токами, а не повреждённые при извлечении численные выражения статьи.

\subsection{Шаг 10. Исключить короткозамкнутые витки в окружающей меди}
Проверяют все медные полигоны, экраны, крепёжные пластины и рамки, которые способны образовать замкнутый путь вокруг потока. Особенно опасны автоматически залитый полигон вокруг выреза и медная рамка, замкнутая через переходы на другом слое. Проверяют не только каждый слой, но и полный трёхмерный граф проводящих соединений.

Электростатический экран выполняют с разрывом, не допускающим замкнутого витка; его подключение выбирают по схеме синфазных токов. Разрыв на одном слое не помогает, если он шунтирован соединением на другом. Проверяют изоляцию экрана от обеих обмоток и вихревые потери по \eqref{eq:metal-eddy-loss}. Неиспользуемую медь в области сильного поля удаляют либо обосновывают её наличие.

\subsection{Шаг 11. Установить правила CAD и проверить выпускные данные}
Минимальный комплект правил включает ширину меди, междорожечные зазоры, внутрицепные зазоры витков, отверстие--медь, отверстие--отверстие, медь--вырез, отступ от зазора сердечника, первично-вторичный барьер и ограничения каждого перехода. В правилах отдельно различают технологическую и электрическую изоляцию.

После трассировки:
\begin{enumerate}
\item выполнить DRC и независимую проверку схемы последовательных и параллельных соединений;
\item экспортировать Gerber X2 либо согласованный ODB++, файлы сверления PTH/NPTH, фрезерование, стек и чертёж;
\item открыть экспорт в независимом CAM-просмотрщике, проверить порядок слоёв, зеркальность, единицы и совмещение отверстий;
\item повторно проследить медный путь, разрывы витков и изоляционные расстояния уже в экспортированных данных;
\item получить от изготовителя DFM-отчёт и перечень предлагаемых CAM-изменений; изменения критического рисунка утверждать явно.
\end{enumerate}
\source Рекомендации по слоям, полигонам, узким остаткам меди и документации на порядок слоёв --- \cite{web17}; правила сверления --- \cite{web18}.

\subsection{Учебный пример выбора ширины, не проект конкретного изделия}
\stepgoal{Показать расчётный переход от доступного окна к CAD-ширине и проверить недостаточность одной DC-оценки.}
Предполагается полоса шириной 8,0 мм, четыре соседних участка витков, проектный междорожечный зазор 0,30 мм и краевые резервы по 0,50 мм. Допуски и ток заданы \emph{только для иллюстрации}; они не являются нормами завода. Принимаются CAD-ширина 1,50 мм, возможное уменьшение ширины 0,10 мм, минимальная готовая медь 60 мкм и длина одной полной ветви 0,20 м. Предусмотрены две одинаковые параллельные ветви, общий постоянный ток 4 А. Межслойные переходы в числовую оценку пока не включены.

По условию размещения \eqref{pl:eq:radial-fit} предельная ширина составляет $1{,}525$ мм; выбранные $1{,}50$ мм занимают вместе с зазорами и резервами $7{,}90$ мм. Это проверка заданных проектных контуров; если у готовой меди возможен положительный допуск, его проверяют отдельно.

По \eqref{pl:eq:minimum-section} минимальная ширина равна $1{,}40$ мм, минимальное сечение каждой ветви --- $0{,}084$ мм$^2$. При равномерном делении ток каждой ветви равен $2$ А; диагностическая плотность тока составляет $23{,}81$ А/мм$^2$. Это высокая величина, требующая теплового обоснования, а не выбранная допустимая норма.

При удельном сопротивлении меди с принятыми номинальными параметрами и применении \eqref{eq:planar-rdc} получаются сопротивление одной ветви $41{,}05$ мОм и двух ветвей в параллели $20{,}52$ мОм при $20\,^{\circ}\mathrm C$. Постоянные медные потери, рассчитанные по \eqref{eq:copper-loss-harmonic} при отсутствии гармоник, равны $0{,}328$ Вт.

При \emph{заданной для проверки}, а не рассчитанной температуре $100\,^{\circ}\mathrm C$ по \eqref{eq:rho-temperature} сопротивление двух ветвей составляет $26{,}98$ мОм, потери --- $0{,}432$ Вт. Эта подстановка не означает, что установившаяся температура будет $100\,^{\circ}\mathrm C$.

Для дополнительного понимания масштаба: по \eqref{eq:skin-depth} на $150$ кГц при $20\,^{\circ}\mathrm C$ глубина проникновения равна $170{,}6$ мкм. Даже если толщина меди меньше этой величины, эффект близости и поле зазора не исключены.


\criterion Даже геометрически помещающаяся дорожка может иметь неприемлемые потери. Пример не подтверждает допустимый ток платы: необходимо добавить переходы, переменную составляющую, токораспределение и тепловую модель. Если допустимый нагрев не обеспечен, изменяют ширину, число параллельных ветвей, стек, размер сердечника либо технологию; не объявляют выбранную плотность тока нормативной.

\section{РАСЧЁТ ПОТЕРЬ И ВЫБОР ЧЕРЕДОВАНИЯ ОБМОТОК}
\label{pl:losses}
\subsection{Постоянный ток и реальная геометрия}
Используют \eqref{eq:rho-temperature}, \eqref{eq:planar-rdc} и \eqref{eq:via-resistance}. Сопротивления считают по участкам, включая выводы и сужения. Если параллельные пути различны, сначала получают отдельные сопротивления и решают электрическую сеть; сокращённое деление на число параллельных слоёв не применяют без проверки.

\subsection{Скин-эффект, эффект близости и краевое поле}

\begin{equation}
 \delta(f,\vartheta)=\sqrt{\frac{\rho(\vartheta)}{\pi f\mu_{\rm Cu}}}.
 \label{eq:skin-depth}
\end{equation}
$\delta$ --- глубина проникновения, м; $f>0$ --- частота гармоники, Гц; для меди $\mu_{\rm Cu}\simeq\mu_0\simeq1{,}256637\cdot10^{-6}$ Гн/м. Это решение однородной проводящей среды, не формула полного сопротивления печатной спирали.
\begin{equation}
 P_{\rm Cu}=I_{\rm DC}^2R_{\rm DC}+\sum_{h\geq1}I_{h,\rm rms}^2R_{\rm AC}(f_h).
 \label{eq:copper-loss-harmonic}
\end{equation}
Сумма справедлива для линейной модели меди при фиксированных температуре и условиях возбуждения. $R_{\rm AC}$ учитывает скин-эффект, близость и внешнее поле; его нельзя брать из холостого хода трансформатора для произвольной нагрузки. Для взаимодействующих ветвей используют матричную форму ниже.
\begin{equation}
 P_{\rm metal,h}=\int_{V_{\rm metal}}\frac{|\mathbf J_{h,\rm rms}|^2}{\sigma}\,dV.
 \label{eq:metal-eddy-loss}
\end{equation}
$\sigma$ --- проводимость, См/м, $\mathbf J_{h,\rm rms}$ --- комплексная действующая плотность тока. При пиковых амплитудах в интеграле требуется множитель $1/2$. Конвенцию экспортируемого FEM-поля фиксируют в отчёте.

Глубина проникновения определяется \eqref{eq:skin-depth}. Сравнение толщины меди с глубиной проникновения служит первичным фильтром, но не критерием отсутствия AC-потерь: поле соседней обмотки и поперечное поле зазора могут вытеснять ток по ширине даже тонкой дорожки. Для дросселя с DC-смещением DC-потери и гармонические потери разделяют согласно \eqref{eq:copper-loss-harmonic}.\cite{web2,schaefer,ewaldfoil}

Одномерную модель Доуэлла допустимо использовать для быстрого сравнения слоёв, если поле приблизительно одномерно, проводники близки к сплошным слоям и нет доминирующего краевого поля зазора. Её не распространяют без проверки на частично заполненную спираль, площадки переходов и близкую к зазору широкую медь. В итоговой модели учитывают фактический рисунок и влияние экранирования на индуктивность.

\subsection{Почему одного $R_{\rm AC}$ трансформатора недостаточно}
При одной и той же частоте потери первичной меди зависят от токов соседних обмоток. Значение, полученное при разомкнутой вторичной обмотке, не является потерями под нагрузкой.

Для фиксированной температуры и линейной электромагнитной модели проводников медные потери гармоники можно записать как квадратичную форму
\begin{equation}
 P_{\mathrm{Cu},h}=\mathbf I_h^{\dagger}\,\mathbf R_{\mathrm{Cu},h}\,\mathbf I_h .
 \label{pl:eq:loss-matrix}
\end{equation}
Здесь $\mathbf I_h$ --- вектор комплексных \emph{действующих} токов независимых ветвей, $\dagger$ --- эрмитово сопряжение, $\mathbf R_{\mathrm{Cu},h}$ --- эрмитова положительно полуопределённая матрица именно медных потерь. Для действующих амплитуд дополнительный множитель $1/2$ не нужен. Матрицу получают из полевой модели с независимыми возбуждениями и перекрёстными членами либо вычисляют потери непосредственно при рабочем наборе токов. Действительная часть полного входного импеданса может содержать потери сердечника и не должна автоматически приниматься за медные потери.

Это обобщение скалярного расчёта для нескольких взаимодействующих проводников. При существенно нелинейном сердечнике гармонические постановки не заменяют восстановление рабочей траектории намагничивания и отдельный расчёт потерь сердечника.

\subsection{Сравнение вариантов interleaving}
\begin{longtable}{>{\raggedright\arraybackslash}p{30mm}>{\raggedright\arraybackslash}p{62mm}>{\raggedright\arraybackslash}p{62mm}}
\caption{Кандидаты на расположение секций}\label{pl:interleave}\\
\toprule Схема & Ожидаемая тенденция & Обязательная обратная проверка\\\midrule\endfirsthead
\toprule Схема & Ожидаемая тенденция & Обязательная обратная проверка\\\midrule\endhead
P--P--S--S & Меньше интерфейсов P--S, проще разделить барьер & Возможны большие рассеяние и эффект близости\\
P/2--S--P/2 & Компенсация ампер-витков по окну, потенциально меньше рассеяние & Два интерфейса барьера, ёмкость, переходы секций P\\
P--S--P--S & Более сильное чередование & Много интерфейсов P--S, сложнее изоляция и CM-пути\\
Раздельные PCB & Гибкость изоляционного расстояния и сборки & Соединения, увеличенное рассеяние, высота пакета\\
\bottomrule
\end{longtable}
Тенденции не являются ранжированием всех реальных конструкций. Для каждой схемы задают полные ампер-витки слоя и токи преобразователя: близость геометрических слоёв сама по себе не гарантирует компенсацию магнитного поля. Сравнивают при одинаковых витках, материале, условиях охлаждения и требованиях изоляции.\cite{web1,web2}

\criterion Выбирают не максимальное чередование, а допустимый компромисс потерь меди, рассеяния, ёмкости, технологичности и стоимости. Компенсирующие экраны и дополнительные переходы включают в те же расчёты, а не считают бесплатным улучшением.

\section{ПАРАЗИТНЫЕ ПАРАМЕТРЫ И ЭЛЕКТРОМАГНИТНАЯ СОВМЕСТИМОСТЬ}
\label{pl:parasitics}
\subsection{Потенциалы слоёв и межобмоточная ёмкость}

\begin{equation}
 W_e=\frac12 C_{\rm eq}U^2=\frac12\int_V\mathbf E\cdot\mathbf D\,dV,
 \qquad C_{\rm eq}=\frac{2W_e}{U^2}.
 \label{eq:capacitance-energy}
\end{equation}
$U$ --- статическое испытательное напряжение, В; $W_e$ --- энергия, Дж. Для получения двухполюсной ёмкости фиксируют потенциалы всех остальных проводников либо условия плавающих электродов. Общая многополюсная форма $W_e=\mathbf v^T\mathbf C\mathbf v/2$ требует одной согласованной системы отсчёта.
\begin{equation}
 i_{\rm CM}=C_{\rm CM}\frac{du_{\rm CM}}{dt},\qquad
 I_{{\rm CM},h}=2\pi f_h C_{\rm CM}|U_{{\rm CM},h}|.
 \label{eq:common-mode-current}
\end{equation}
$C_{\rm CM}$ --- ёмкость выбранного пути; гармонические напряжения и токи --- действующие. При нескольких путях токи складывают комплексно, с их фазами.
\begin{equation}
 Z(\omega)=\left[\frac1{R+j\omega L}+j\omega C\right]^{-1},\qquad
 f_{\rm SR}\simeq\frac1{2\pi\sqrt{LC}}.
 \label{eq:impedance}
\end{equation}
Это однорезонансная модель последовательной ветви $R,L$, шунтированной $C$; оценка частоты предполагает малые потери. Многообмоточная плата требует проверки нескольких мод и условий замыкания выводов.
\begin{equation}
 P_{\rm diel}=\sum_h 2\pi f_h C_h U_{h,\rm rms}^2\tan\delta_{d,h}.
 \label{eq:dielectric-loss-independent}
\end{equation}
$\tan\delta_d$ --- тангенс диэлектрических потерь; здесь он не обозначает глубину проникновения меди. Для связанной системы выражение применяют к согласованным ветвям ёмкостной модели, не суммируя одну и ту же энергию дважды.

Для приблизительно плоского однородного зазора между перекрывающимися проводниками:
\begin{equation}
 C_{ij}^{(0)}=\varepsilon_0\varepsilon_{r,ij}\frac{A_{ij}}{d_{ij}}.
 \label{pl:eq:plate-capacitance}
\end{equation}
Здесь $A_{ij}$ --- площадь перекрытия, м$^2$; $d_{ij}$ --- толщина диэлектрика, м. Это модель параллельных пластин без краевых полей, полученная из однородного электростатического поля. Она является оценкой отдельного геометрического интерфейса, а не автоматически клеммной ёмкостью трансформатора.

Для перехода к эквивалентной ёмкости используют энергетический порядок из \eqref{eq:capacitance-energy} и \eqref{eq:capacitance-energy}, задавая реальные потенциалы витков и экранов. Не складывают все площадные ёмкости арифметически без схемы соединения. Краевые электрические поля и неоднородность диэлектрика учитывают численно, если простая модель не обеспечивает точности.\cite{zhaocap}

\subsection{Синфазные токи, экран и резонансы}
Синфазный ток оценивают по \eqref{eq:common-mode-current} и \eqref{eq:common-mode-current}; воздействующее напряжение --- между реальными переключающимися узлами, а не только постоянное входное напряжение. Проверяют связь обеих обмоток с сердечником, экраном, радиатором и корпусом.

Экран способен изменить маршрут токов смещения, но его эффективность зависит от точки подключения и её импеданса. Подключение экрана к произвольной «земле» может увеличить помеху. Широкий экран проверяют на вихревые токи и сохранение изоляции; принцип разомкнутого витка из раздела \ref{pl:pcb} обязателен.

Однорезонансные выражения \eqref{eq:impedance} и \eqref{eq:impedance} используют лишь там, где соответствующая сосредоточенная схема подтверждена. Многообмоточный планарный компонент может иметь несколько мод; измеряют частотную характеристику при разных состояниях остальных выводов. Уменьшение измеренной ёмкости в одной схеме подключения не доказывает уменьшения CM-тока в преобразователе.

\section{ИЗОЛЯЦИЯ, МАТЕРИАЛЫ ПЛАТЫ И ТЕХНОЛОГИЯ ИЗГОТОВЛЕНИЯ}
\label{pl:manufacturing}
\subsection{Порядок выбора изоляции}
\begin{enumerate}
\item Определить применимый стандарт изделия и категорию изоляции; выбрать контрольные пары цепей P--S, P--core, S--core, экран--обмотки и крепёж--обмотки.
\item Установить рабочее напряжение, повторяющиеся пики, импульсные воздействия, условия среды, высоту, загрязнение и требования к материалу.
\item По применимым нормам определить воздушные расстояния, пути утечки, требования к твёрдой изоляции и испытаниям.
\item Проверить все пути: поверхность платы, вырез под сердечник, торцы, стенки отверстий, выводы, клей, крепёж и радиатор.
\item Подтвердить фактический стек и изоляцию испытаниями и микрошлифами, если это предусмотрено программой качества.
\end{enumerate}
\source Термины, различие типов изоляции и связь со стандартами приведены в ANP151.\cite{web25} Полные нормативные тексты IPC и IEC в настоящей работе не были получены; численные нормативные таблицы по ним не воспроизводятся. Это не декларация соответствия серийного изделия стандарту.

\criterion Паяльная маска не считается автоматически основной или усиленной изоляцией. Паспортная кратковременная электрическая прочность ламината не равна допустимому длительному рабочему напряжению PCB. Препрег, клей, лак и заливка учитываются как барьер только после квалификации соответствующей технологии. При пониженном давлении и для космического применения отдельно рассматривают разряды, дегазацию, термоциклы и применимые отраслевые требования.

\subsection{Изготовление рисунка и сборка многослойной платы}
Типовой маршрут вычитающей технологии включает подготовку медного ламината, формирование защитного рисунка, травление внутренних слоёв, оптический контроль, укладку и прессование пакета, сверление, подготовку стенок отверстий, металлизацию, формирование наружного рисунка, маску, финишное покрытие, фрезерование и электрический контроль. Реальная последовательность наружного травления и гальваники зависит от технологии завода.\cite{web16,web17,web18,web22}

\method В технических требованиях нормируют не только «толстую медь», но и минимальные готовые ширину и толщину в критических участках, остаточный диэлектрик, металлизацию отверстий, допустимый уход слоёв и качество кромок. Для толстой меди согласуют заполнение смолой, отсутствие пустот и требования к плоскостности.

Не допускают без согласования: уменьшение критического отверстия, удаление переходов, добавление медных перемычек, заливку технологических пустот металлом, перенос выреза и уменьшение барьера. Балансирующую медь разрешают только вне запрещённых полевых и изоляционных зон либо после отдельной проверки.

\subsection{Гибкая и гибко-жёсткая печатная обмотка}
Если трансформатор требует значительно больше слоёв, чем остальная схема, сравнивают отдельную многослойную плату и складываемую гибкую обмотку. Последняя позволяет не распространять число обмоточных слоёв на всю площадь преобразователя, но экономию рассчитывают по реальным расценкам, раскрою и выходу годных. Оценочные коэффициенты стоимости из старой публикации не являются современной ценой производства.\cite[раздел II]{doi54}

При проектировании складываемой обмотки выпускают развёртку, порядок сгибов, полярность каждого участка и контрольную модель сложенного пакета. Учитывают конечный радиус изгиба и наложение сгибов: сумма толщин плоских слоёв может заметно недооценить высоту у кромки. Схемы сгиба по одной или двум осям сравнивают по использованию материала, высоте и доступности сборки.

\criterion В экспериментальной работе de Jong--Ferreira--Bauer отмечены микротрещины меди на сгибах и отрыв соединений у microvia вблизи сгиба. Поэтому радиус, материал меди, число допустимых сгибаний, отступ переходов и электрический контроль после формовки задают с изготовителем. Утверждение авторов об однократном складывании не заменяет квалификацию термоциклов и вибрации. Проверяют сопротивление и изоляцию до и после формовки и климатических испытаний.\cite[раздел IV, с.~3162]{doi54}

\subsection{Сборка сердечника и фиксация зазора}
\begin{enumerate}
\item Проверить платы до установки сердечника: размеры, отсутствие замыкания витков, сопротивление и изоляционные расстояния.
\item Осмотреть феррит на сколы, проверить сопрягаемые поверхности и размеры зазора.
\item Установить квалифицированные прокладки и средства фиксации; не считать неконтролируемый клеевой шов точным магнитным зазором.
\item Собрать половины с установленным усилием крепления без деформации PCB; проверить отсутствие контакта кромки феррита с медью.
\item После отверждения повторно измерить индуктивность и сопротивление, затем выполнить предусмотренные испытания изоляции.
\end{enumerate}
\recommend Шлифование феррита и изготовление прецизионного зазора поручают квалифицированному исполнителю с контролем размеров и пыли. Геометрия, достигнутая случайной толщиной клея, не является воспроизводимой серийной конструкцией.

\subsection{Купоны и входной контроль}
Для опытной партии предусматривают купон для микрошлифа стека, минимальной дорожки, наиболее нагруженных переходов и изоляционного барьера. Купон должен соответствовать критическим участкам по технологии и условиям металлизации; произвольный образец с края панели не всегда представляет худшую область.

Контролируют готовую медь, подтрав, остаточный диэлектрик, совмещение, металлизацию стенки, пустоты, трещины и чистоту после сборки. Измерение общей толщины платы не подтверждает каждую межслойную толщину. Отчёт производителя и фактическая ревизия Gerber связываются с номером партии.

\section{ВЕРИФИКАЦИЯ МЕТОДОМ КОНЕЧНЫХ ЭЛЕМЕНТОВ}
\label{pl:fem}
\subsection{Последовательность постановок}
Расчётная геометрия должна содержать все существенные медные слои, переходы и выводы. Рекомендуемая последовательность:
\begin{enumerate}
\item магнитостатическая либо нелинейная проверка рабочего потока и запаса по насыщению;
\item вихретоковая постановка по существенным гармоникам с реальными токами всех обмоток;
\item электростатическая постановка с действительными потенциалами и экранами;
\item тепловая постановка с пространственным распределением полученных потерь;
\item повторение электрической и тепловой частей при обновлённой температуре материалов.
\end{enumerate}
Нелинейное намагничивание и линейная AC-модель меди не смешиваются без явного описания того, что передано между постановками. Заданная равномерная плотность тока в stranded-модели не доказывает отсутствие вытеснения тока в solid-дорожке.

\subsection{Геометрия, возбуждение и сетка}
Двумерная модель применяется для предварительного отбора сечений. Трёхмерная проверка требуется, когда существенны переходы, углы, выводы, частичное перекрытие слоёв, разрывы экранов или неравные ветви. В токовом порту задают реальный проводящий путь; замкнутый случайный solid и электрически правильная обмотка не эквивалентны.

Уточняют сетку по толщине меди, глубине проникновения, кромкам и зазору. Если используется импедансная граница или тонкослойная модель, проверяют её применимость к толщине и локальному полю. Верификация проводится по интегральным потерям и параметрам, а не только по красивой карте $J$.

\subsection{Проверки результата}

\begin{equation}
 \varepsilon_Q=\frac{|Q_m-Q_{m-1}|}{\max(|Q_m|,Q_{\rm floor})}.
 \label{eq:mesh-convergence}
\end{equation}
$Q_m$ --- контролируемый результат на сетке $m$; $Q_{\rm floor}>0$ --- обоснованный масштаб для малых величин. Сравнивают не менее трёх сопоставимых сгущений без изменения физической постановки; отдельно проверяют удаление внешней границы. Допустимую ошибку задают из требований проекта.
\begin{equation}
 W_{\rm field}\simeq\frac12\mathbf i^T\mathbf L\mathbf i
 \quad\text{(для одной обмотки: }W_{\rm field}\simeq LI^2/2\text{)}.
 \label{eq:fem-energy-check}
\end{equation}
Это контроль линейной магнитостатической модели с мгновенными токами. Он не распространяется без изменения на нелинейную гистерезисную модель либо комплексную гармоническую энергию.

Применяют \eqref{eq:mesh-convergence} и \eqref{eq:fem-energy-check}. Дополнительно проверяют:
\begin{itemize}
\item непрерывность тока, сумму токов ветвей и отсутствие неучтённых циркулирующих контуров;
\item согласованность собственных и взаимных индуктивностей с выбранной схемой, взаимность и пассивность линейной матрицы;
\item раздельные потери меди, сердечника, экранов и прочих проводящих деталей;
\item чувствительность к положению границ области, толщине меди, смещению слоёв и положению зазора;
\item сходимость температуры после передачи потерь в тепловую постановку.
\end{itemize}
\criterion Требуемую точность назначают из допустимой погрешности проекта. Несколько адаптивных проходов сами по себе не являются доказательством сходимости. Пиковое поле на идеальном остром ребре не используют как критерий разрушения или насыщения без анализа геометрии и сетки.

\section{ТЕПЛОВОЙ РАСЧЁТ И СОВМЕСТНАЯ ИТЕРАЦИЯ}
\label{pl:thermal}

\begin{equation}
 q_{v,k}=P_k/V_k,\qquad \Delta\vartheta=P_\Sigma R_\theta.
 \label{eq:heat-density}
\end{equation}
$q_{v,k}$ --- средняя плотность тепловыделения, Вт/м$^3$; $R_\theta$ --- тепловое сопротивление при заданных условиях, К/Вт. Усреднение источника не описывает локальную горячую точку. Для уточнения решают
\begin{equation}
 \rho_m c_p\frac{\partial\vartheta}{\partial t}
 =\nabla\cdot(\mathbf K\nabla\vartheta)+q_v.
 \label{eq:thermal-pde}
\end{equation}
$\rho_m$ --- массовая плотность, кг/м$^3$ (не электрическое сопротивление); $c_p$ --- Дж/(кг$\cdot$К); $\mathbf K$ --- тензор теплопроводности, Вт/(м$\cdot$К). Задают контактные сопротивления, конвекцию и излучение; в стационарной задаче производная равна нулю. Для вакуума не подставляют земную конвекцию.

Плотность источников тепла и начальную тепловую оценку определяют по \eqref{eq:heat-density} и \eqref{eq:heat-density}. Для PCB одной общей тепловой постоянной обычно недостаточно: теплопроводность вдоль меди и поперёк диэлектрика различается, а внутренние слои и переходы могут быть горячее поверхности.

\method В модель включают ламинат, медь каждого слоя либо проверенную эквивалентную анизотропную среду, сердечник, прокладки, клей, радиатор и реальные условия теплообмена. Тепловые переходы через изоляционный барьер не размещают без электрической проверки. Контакт с радиатором учитывает и тепловое сопротивление, и ёмкость к корпусу.

\begin{enumerate}
\item Рассчитать потери при начальной температуре.
\item Получить температуру дорожек, переходов, диэлектрика, сердечника и соединений.
\item Обновить сопротивление меди, свойства феррита и диэлектрика.
\item Повторять до согласованного изменения температур и потерь.
\item Проверить худшие сочетания температуры окружающей среды, нагрузки, частоты и допусков.
\end{enumerate}
\criterion Температуру ограничивает наиболее слабый материал или соединение, а не только медь. Температура стеклования не является автоматически допустимой длительной рабочей температурой. В отчёте нельзя выдавать оценку по $R_\theta$ за результат верифицированной тепловой модели.

\section{ОПТИМИЗАЦИЯ И КРИТЕРИИ ПРИЁМКИ}
\label{pl:acceptance}
\subsection{Переменные и сопоставление вариантов}

\begin{equation}
 \overline P_\Sigma=\sum_r w_rP_{\Sigma,r},\qquad
 w_r\geq0,\quad\sum_r w_r=1.
 \label{eq:mission-profile-loss}
\end{equation}
$w_r$ --- доля времени режима. Оптимизация среднего не отменяет ограничений в наиболее тяжёлом режиме. $P_\Sigma$ включает медь, сердечник, диэлектрик, экраны и соединения без двойного счёта. При цикле, сравнимом с тепловыми постоянными, рассчитывают переходную температуру, а не только среднюю мощность.

Варьируют материал, типоразмер, число витков, зазор, порядок слоёв, число последовательных секций и параллельных ветвей, готовую толщину меди, ширину рисунка, число и расположение переходов, межслойные толщины, форму выводов и положение зазора относительно PCB. Оптимизация ведётся по полному профилю режимов из \eqref{eq:mission-profile-loss}.

Сначала строят карту режимов фиксированной конструкции. Затем изменяют по одному обоснованному набору параметров и сравнивают с тем же профилем. Иначе улучшение, полученное одновременно изменением режима преобразователя и конструкции, нельзя приписывать только планарной обмотке.

Для каждого кандидата фиксируют потери, максимальные температуры, индуктивности, токи ветвей, ёмкости, запас изоляции, габариты, стоимость и технологические риски. Сохраняют отвергнутые варианты с причиной отказа. Предпочтение минимальной стоимости не отменяет обязательных ограничений.

\subsection{Испытания опытного образца}
\begin{longtable}{>{\raggedright\arraybackslash}p{35mm}>{\raggedright\arraybackslash}p{69mm}>{\raggedright\arraybackslash}p{50mm}}
\caption{Программа верификации}\label{pl:tests}\\
\toprule Проверка & Условия и измерение & Что подтверждается\\\midrule\endfirsthead
\toprule Проверка & Условия и измерение & Что подтверждается\\\midrule\endhead
Голая PCB & CAM-связность, AOI, электрический тест, микрошлиф по программе & Рисунок, стек, переходы и технологические минимумы\\
Фазировка и витки & Малосигнальный источник; отдельно все выводы и ветви & Полярность и отношение ЭДС без силового риска\\
DC-сопротивление & Четырёхпроводное измерение, известная температура & Медь и соединения; сравнение с расчётом\\
Дроссель & $L(I,T)$, пульсация, импеданс; уточнить дифференциальную величину & Рабочая индуктивность и нелинейность\\
Трансформатор & Open/short, коэффициент передачи; заданные частота и амплитуда & Согласованность эквивалентной схемы, намагничивания и рассеяния\\
Ёмкость и резонансы & Разные соединения остальных выводов; компенсация оснастки & Многополюсные паразитные параметры\\
Рабочие потери & Согласованные формы напряжений и токов; оценка фазовой погрешности & Потери в реальном возбуждении\\
Нагрев & Номинальный и худшие режимы, установившееся состояние & Температуры с учётом погрешности и недоступных внутренних слоёв\\
Изоляция & Утверждённые напряжения, время, ток отсечки и пары цепей & Только предусмотренные программой критерии\\
Надёжность & Термоциклы, механика, влажность и повторные измерения по ТЗ & Устойчивость платы и сборки к среде\\
\bottomrule
\end{longtable}
Измерения высоковольтной изоляции выполняют на предназначенном для этого оборудовании и по утверждённой безопасной процедуре. Универсального испытательного напряжения для всех планарных трансформаторов нет.

\subsection{Проверка численного и экспериментального согласия}
В отчёте для каждого сравнения приводят частоту, амплитуду, температуру, состояние остальных обмоток и оснастку. Расхождение $L$, измеренной при малом сигнале без DC, и рабочей дифференциальной индуктивности под нагрузкой не объявляют ошибкой FEM до согласования условий.

Потери малого высокодобротного компонента чувствительны к фазовой ошибке измерения; при необходимости используют калориметрию либо независимую проверку. Тепловизор измеряет доступную поверхность с зависимостью от излучательной способности и не показывает непосредственно температуру внутренней меди.

\criterion Приёмка разрешается только после закрытия электрических, тепловых, изоляционных и технологических требований. Настоящий универсальный документ не задаёт численных пределов вместо ТЗ и не является протоколом испытанного изделия.

\section{СОСТАВ ИТОГОВОЙ РАСЧЁТНОЙ ДОКУМЕНТАЦИИ}
\label{pl:docs}
Итоговый комплект должен включать:
\begin{enumerate}
\item исходные данные, режимы, применимые стандарты и матрицу требований;
\item принятые физические модели, область их применимости, систему обозначений и версию расчёта;
\item расчёт витков, материала и зазора; схему эквивалентного компонента;
\item электрическую таблицу слоёв, секций, полярностей и параллельных ветвей;
\item исходный CAD-проект, Gerber/ODB++, PTH/NPTH, фрезерование, стек и сборочный чертёж;
\item расчёт минимального сечения, переходов, DC- и AC-потерь, ёмкостей и нагрева;
\item FEM-модели, свойства материалов, возбуждения, журналы сходимости и карты $B$, $H$, $J$ с единицами; температуры только при выполненном тепловом расчёте;
\item DFM-отчёт и утверждённые изменения изготовителя; паспорт партии и результаты купонов;
\item программу испытаний, сырые измерения, погрешности, сравнение с расчётом и решение о пригодности;
\item реестр источников с указанием реально доступных полных текстов и ограничений доказательной базы.
\end{enumerate}

\subsection{Содержание чертежа печатной обмотки}
На чертеже указывают начало и конец каждой обмотки, число витков, связи слоёв, готовую медь, минимальные ширины и зазоры, изоляционные области, размеры отверстий и переходов, ограничения по изменениям CAM, положение зазора сердечника, прокладки и крепёж. В таблице исполнения фиксируют ревизию рисунка, материал платы, магнитопровод, материал феррита и зазор.

\recommend Комплект связывают единым идентификатором исполнения и ревизии. Контрольные суммы CAD, Gerber, расчётных данных и протоколов позволяют проверить, что испытывался именно рассчитанный рисунок. Любое изменение стека, переходов, изоляции или зазора требует оценки влияния и повторных проверок.

\section{ЗАКЛЮЧЕНИЕ}
Расчёт планарного дросселя объединяет электрический режим, магнитную цепь, геометрию PCB, токораспределение, краевое поле, нагрев и производство. Для трансформатора дополнительно обязательны многообмоточная модель, выбор коэффициента трансформации, проверка намагничивания и рассеяния, фазировка и изоляционный барьер.

Печатная дорожка проектируется в замкнутой последовательности: электрическая карта витков --- механические ограничения --- минимальное готовое сечение --- размещение --- переходы и стек --- потери и нагрев --- DFM --- изготовление и измерение. Ни выбранная плотность тока, ни толщина меди, ни отсутствие ошибок обычного DRC по отдельности не доказывают пригодность обмотки.

\appendix
\section{СОСТОЯНИЕ ПРОВЕРКИ ЛИТЕРАТУРЫ}
Список включает исходную подборку из 89 позиций. Она сохранена для прослеживаемости, а не выдана за 89 полностью прочитанных работ. По каждому источнику ведётся отдельный журнал получения текста и чтения. Получение PDF, извлечение текста и содержательная проверка --- разные этапы.

Материалы, полный текст которых пока не получен, не являются подтверждением заимствованных формул. Такие позиции отмечены в списке. Для сложных моделей без законченной проверки первоисточника обязательна самостоятельная проверка нормировки, пределов применимости и сопоставление с измерениями. Формулы общего электромагнитного расчёта изложены непосредственно и не требуют другого документа.

Подробный журнал находится в сопровождающем файле \nolinkurl{SOURCE_REVIEW_STATUS.md}. До закрытия его незавершённых пунктов литературная верификация настоящей редакции считается незавершённой.

\begin{thebibliography}{99}
\small\raggedright
\bibitem[1]{web1}
\textit{Texas Instruments. Magnetics Design 4 - Power Transformer Design}. .
\url{https://www.ti.com/lit/ml/slup126/slup126.pdf}.
\textit{Доступность:} Полный PDF получен; доступность не означает постраничного рецензирования. \textit{Тема:} Апнот / руководство / производство.
\bibitem[2]{web2}
\textit{Texas Instruments. Magnetics Design 3 - Windings}. .
\url{https://www.ti.com/lit/ml/slup125/slup125.pdf}.
\textit{Доступность:} Полный PDF получен; доступность не означает постраничного рецензирования. \textit{Тема:} Апнот / руководство / производство.
\bibitem[3]{web3}
\textit{Texas Instruments. Magnetics Design 2 - Magnetic Core Characteristics}. .
\url{https://www.ti.com/lit/ml/slup124/slup124.pdf}.
\textit{Доступность:} Полный PDF получен; доступность не означает постраничного рецензирования. \textit{Тема:} Апнот / руководство / производство.
\bibitem[4]{web35}
\textit{Texas Instruments. Magnetics Design 5 - Inductor and Flyback Transformer Design}. .
\url{https://www.ti.com/lit/ml/slup127/slup127.pdf}.
\textit{Доступность:} Полный PDF получен; доступность не означает постраничного рецензирования. \textit{Тема:} Апнот / руководство / производство.
\bibitem[5]{web25}
\textit{Würth Elektronik. ANP151a — Mastering Galvanic Isolation: Ensuring Safety in High-Power Electronics}. .
\url{https://www.we-online.com/components/media/o896744v410%20ANP151_EN.pdf}.
\textit{Доступность:} Полный PDF получен; доступность не означает постраничного рецензирования. \textit{Тема:} Апнот / руководство / производство.
\bibitem[6]{web17}
\textit{Copper Layers - Eurocircuits}. .
\url{https://www.eurocircuits.com/technical-guidelines/pcb-design-guidelines/copper-layers/}.
\textit{Доступность:} Веб-страница получена. \textit{Тема:} Апнот / руководство / производство.
\bibitem[7]{web18}
\textit{Drilled Holes - Eurocircuits}. .
\url{https://www.eurocircuits.com/technical-guidelines/pcb-design-guidelines/drilled-holes/}.
\textit{Доступность:} Веб-страница получена. \textit{Тема:} Апнот / руководство / производство.
\bibitem[8]{web16}
\textit{Track Width and Isolation Gap Tolerances - Eurocircuits}. .
\url{https://www.eurocircuits.com/pcb-design-guidelines/track-width-and-isolation-gap-tolerances/}.
\textit{Доступность:} Веб-страница получена. \textit{Тема:} Апнот / руководство / производство.
\bibitem[9]{web22}
\textit{Tolerances on Copper Thickness - Eurocircuits}. .
\url{https://www.eurocircuits.com/technical-guidelines/understanding-manufacturing-tolerances-on-a-pcb/tolerances-on-copper-thickness/}.
\textit{Доступность:} Веб-страница получена. \textit{Тема:} Апнот / руководство / производство.
\bibitem[10]{web21}
\textit{Layer to Layer Registration Tolerances - Eurocircuits}. .
\url{https://www.eurocircuits.com/technical-guidelines/understanding-manufacturing-tolerances-on-a-pcb/layer-to-layer-registration-tolerances/}.
\textit{Доступность:} Веб-страница получена. \textit{Тема:} Апнот / руководство / производство.
\bibitem[11]{web20}
\textit{Finished PCB Thickness Tolerance - Eurocircuits}. .
\url{https://www.eurocircuits.com/technical-guidelines/understanding-manufacturing-tolerances-on-a-pcb/finished-pcb-thickness-tolerances/}.
\textit{Доступность:} Веб-страница получена. \textit{Тема:} Апнот / руководство / производство.
\bibitem[12]{web19}
\textit{Copper Distribution on a Panel - Eurocircuits}. .
\url{https://www.eurocircuits.com/technical-guidelines/panel-guidelines/copper-distribution-on-a-panel/}.
\textit{Доступность:} Веб-страница получена. \textit{Тема:} Апнот / руководство / производство.
\bibitem[13]{web23}
\textit{Magnetics - Ferrite Core Inductor Design}. .
\url{https://www.mag-inc.com/Design/Design-Guides/Inductor-Design-with-Magnetics-Ferrite-Cores.aspx}.
\textit{Доступность:} Веб-страница получена. \textit{Тема:} Апнот / руководство / производство.
\bibitem[14]{web24}
\textit{Magnetics - Ferrite Core Transformer Design}. .
\url{https://www.mag-inc.com/Design/Design-Guides/Transformer-Design-with-Magnetics-Ferrite-Cores.aspx}.
\textit{Доступность:} Веб-страница получена. \textit{Тема:} Апнот / руководство / производство.
\bibitem[15]{web7}
\textit{WE-PLN Planar Transformers | Passive Components | Würth Elektronik Product Catalog}. .
\url{https://www.we-online.com/en/components/products/WE-PLN}.
\textit{Доступность:} Веб-страница получена. \textit{Тема:} Апнот / руководство / производство.
\bibitem[16]{web26}
\textit{Würth Elektronik. WE-PLN 750341134 — Datasheet}. .
\url{https://www.we-online.com/components/products/datasheet/750341134.pdf}.
\textit{Доступность:} Полный PDF получен; доступность не означает постраничного рецензирования. \textit{Тема:} Апнот / руководство / производство.
\bibitem[17]{web27}
\textit{Products-Ferroxcube}. .
\url{https://www.ferroxcube.com/en-global/products_ferroxcube/detail/shape_cores_accessories/planar_e_cores}.
\textit{Доступность:} Веб-страница получена. \textit{Тема:} Апнот / руководство / производство.
\bibitem[18]{doi18}
I.W. Hofsajer, J.A. Ferreira, D. van Wyk. \textit{Design and analysis of planar integrated L-C-T components for converters}. 2000.
\url{https://doi.org/10.1109/63.892837}.
\textit{Доступность:} Проверена библиографическая запись Crossref; полный текст не получен. \textit{Тема:} Планарные компоненты и схемные модели.
\bibitem[19]{doi19}
M. Meinhardt, M. Duffy, T. O'Donnell, S. O'Reilly, J. Flannery, C. O Mathuna. \textit{New method for integration of resonant inductor and transformer-design, realisation, measurements}. 1999.
\url{https://doi.org/10.1109/apec.1999.750516}.
\textit{Доступность:} Проверена библиографическая запись Crossref; полный текст не получен. \textit{Тема:} Планарные компоненты и схемные модели.
\bibitem[20]{doi20}
J. Biela, J.W. Kolar. \textit{Electromagnetic integration of high power resonant circuits comprising high leakage inductance transformers}. .
\url{https://doi.org/10.1109/pesc.2004.1354802}.
\textit{Доступность:} Проверена библиографическая запись Crossref; полный текст не получен. \textit{Тема:} Планарные компоненты и схемные модели.
\bibitem[21]{doi21}
T. Duerbaum, G. Sauerlaender. \textit{Energy based capacitance model for magnetic devices}. .
\url{https://doi.org/10.1109/apec.2001.911635}.
\textit{Доступность:} Проверена библиографическая запись Crossref; полный текст не получен. \textit{Тема:} Паразитные параметры и ЭМС.
\bibitem[22]{doi22}
N.H. Kutkut, D.M. Divan. \textit{Optimal air-gap design in high-frequency foil windings}. 1998.
\url{https://doi.org/10.1109/63.712317}.
\textit{Доступность:} Проверена библиографическая запись Crossref; полный текст не получен. \textit{Тема:} Потери и зазор.
\bibitem[23]{doi23}
Goran Perica. \textit{Elimination of Leakage Effects Related to the Use of Windings with Fractions of Turns}. 1986.
\url{https://doi.org/10.1109/tpel.1986.4766275}.
\textit{Доступность:} Проверена библиографическая запись Crossref; полный текст не получен. \textit{Тема:} Планарные компоненты и схемные модели.
\bibitem[24]{doi24}
Hai Yan Lu, Jian Guo Zhu, S.Y.R. Hui. \textit{Experimental determination of stray capacitances in high frequency transformers}. 2003.
\url{https://doi.org/10.1109/tpel.2003.816186}.
\textit{Доступность:} Проверена библиографическая запись Crossref; полный текст не получен. \textit{Тема:} Паразитные параметры и ЭМС.
\bibitem[25]{doi25}
J.-P. Vandelac, P.D. Ziogas. \textit{A novel approach for minimizing high-frequency transformer copper losses}. 1988.
\url{https://doi.org/10.1109/63.17944}.
\textit{Доступность:} Проверена библиографическая запись Crossref; полный текст не получен. \textit{Тема:} Потери и зазор.
\bibitem[26]{doi26}
M. Perry. \textit{Multiple Layer Series Connected Winding Design for Minimum Losses}. 1979.
\url{https://doi.org/10.1109/tpas.1979.319520}.
\textit{Доступность:} Проверена библиографическая запись Crossref; полный текст не получен. \textit{Тема:} Потери и зазор.
\bibitem[27]{doi27}
P.L. Dowell. \textit{Effects of eddy currents in transformer windings}. 1966.
\url{https://doi.org/10.1049/piee.1966.0236}.
\textit{Доступность:} Проверена библиографическая запись Crossref; полный текст не получен. \textit{Тема:} Планарные компоненты и схемные модели.
\bibitem[28]{doi28}
Ziwei Ouyang, Ole C. Thomsen, Michael A. E. Andersen. \textit{Optimal Design and Tradeoff Analysis of Planar Transformer in High-Power DC–DC Converters}. 2012.
\url{https://doi.org/10.1109/tie.2010.2046005}.
\textit{Доступность:} Проверена библиографическая запись Crossref; полный текст не получен. \textit{Тема:} Планарные компоненты и схемные модели.
\bibitem[29]{doi29}
W.G. Hurley, E. Gath, J.G. Breslin. \textit{Optimizing the AC resistance of multilayer transformer windings with arbitrary current waveforms}. 2000.
\url{https://doi.org/10.1109/63.838110}.
\textit{Доступность:} Проверена библиографическая запись Crossref; полный текст не получен. \textit{Тема:} Планарные компоненты и схемные модели.
\bibitem[30]{doi30}
J.A. Ferreira. \textit{Improved analytical modeling of conductive losses in magnetic components}. 1994.
\url{https://doi.org/10.1109/63.285503}.
\textit{Доступность:} Проверена библиографическая запись Crossref; полный текст не получен. \textit{Тема:} Потери и зазор.
\bibitem[31]{doi31}
Wei Chen, Yipeng Yan, Yuequan Hu, Qing Lu. \textit{Model and design of PCB parallel winding for planar transformer}. 2003.
\url{https://doi.org/10.1109/tmag.2003.816147}.
\textit{Доступность:} Проверена библиографическая запись Crossref; полный текст не получен. \textit{Тема:} Токораспределение.
\bibitem[32]{doi32}
W.-J. Gu, R. Liu. \textit{A study of volume and weight vs. frequency for high-frequency transformers}. .
\url{https://doi.org/10.1109/pesc.1993.472059}.
\textit{Доступность:} Проверена библиографическая запись Crossref; полный текст не получен. \textit{Тема:} Планарные компоненты и схемные модели.
\bibitem[33]{doi33}
X. Margueron, J.-P. Keradec, A. Besri. \textit{Current Sharing Between Parallel Turns of a Planar Transformer: Prediction and Improvement Using a Circuit Simulation Software}. 2007.
\url{https://doi.org/10.1109/07ias.2007.273}.
\textit{Доступность:} Проверена библиографическая запись Crossref; полный текст не получен. \textit{Тема:} Токораспределение.
\bibitem[34]{doi34}
R. Prieto, J.A. Cobos, O. Garcia, P. Alou, J. Uceda. \textit{Using parallel windings in planar magnetic components}. .
\url{https://doi.org/10.1109/pesc.2001.954423}.
\textit{Доступность:} Проверена библиографическая запись Crossref; полный текст не получен. \textit{Тема:} Токораспределение.
\bibitem[35]{doi35}
R. Prieto, R. Asensi, J.A. Cobos. \textit{Selection of the appropriate winding setup in planar inductors with parallel windings}. 2010.
\url{https://doi.org/10.1109/ecce.2010.5618404}.
\textit{Доступность:} Полный PDF получен из университетского репозитория; текст прочитан с ограничениями. \textit{Тема:} Токораспределение.
\bibitem[36]{doi36}
David J. Perreault, Jingying Hu, Juan M. Rivas, Yehui Han, Olivia Leitermann, Robert C.N. Pilawa-Podgurski, Anthony Sagneri, Charles R. Sullivan. \textit{Opportunities and Challenges in Very High Frequency Power Conversion}. 2009.
\url{https://doi.org/10.1109/apec.2009.4802625}.
\textit{Доступность:} Проверена библиографическая запись Crossref; полный текст не получен. \textit{Тема:} Планарные компоненты и схемные модели.
\bibitem[37]{doi37}
W.G. Hurley, D.J. Wilcox. \textit{Calculation of leakage inductance in transformer windings}. 1994.
\url{https://doi.org/10.1109/63.285502}.
\textit{Доступность:} Проверена библиографическая запись Crossref; полный текст не получен. \textit{Тема:} Планарные компоненты и схемные модели.
\bibitem[38]{doi38}
Ziwei Ouyang, Ole C. Thomsen, Michael A. E. Andersen. \textit{The analysis and comparison of leakage inductance in different winding arrangements for planar transformer}. 2009.
\url{https://doi.org/10.1109/peds.2009.5385844}.
\textit{Доступность:} Проверена библиографическая запись Crossref; полный текст не получен. \textit{Тема:} Планарные компоненты и схемные модели.
\bibitem[39]{doi39}
B.W. Carsten. \textit{The low leakage inductance of planar transformers; fact or myth?}. .
\url{https://doi.org/10.1109/apec.2001.912515}.
\textit{Доступность:} Проверена библиографическая запись Crossref; полный текст не получен. \textit{Тема:} Планарные компоненты и схемные модели.
\bibitem[40]{doi40}
Jun Zhang, Ziwei Ouyang, M. C. Duffy, M. A. E Andersen, W. G. Hurley. \textit{Leakage inductance calculation for planar transformers with a magnetic shunt}. 2013.
\url{https://doi.org/10.1109/ecce.2013.6646762}.
\textit{Доступность:} Проверена библиографическая запись Crossref; полный текст не получен. \textit{Тема:} Планарные компоненты и схемные модели.
\bibitem[41]{doi41}
K. Kawabe, H. Koyama, K. Shirae. \textit{Planar inductor}. 1984.
\url{https://doi.org/10.1109/tmag.1984.1063271}.
\textit{Доступность:} Проверена библиографическая запись Crossref; полный текст не получен. \textit{Тема:} Планарные компоненты и схемные модели.
\bibitem[42]{doi42}
J.-P. Schauwers, C. Nunes, B. Velaerts, F. Robert, P. Mathys. \textit{Planar transformer technology applied to AC/DC conversion}. 1999.
\url{https://doi.org/10.1109/intlec.1999.794113}.
\textit{Доступность:} Проверена библиографическая запись Crossref; полный текст не получен. \textit{Тема:} Планарные компоненты и схемные модели.
\bibitem[43]{doi43}
W. G. Hurley, W. H. Wölfle. \textit{Transformers and Inductors for Power Electronics}. 2013.
\url{https://doi.org/10.1002/9781118544648}.
\textit{Доступность:} Проверена библиографическая запись Crossref; полный текст не получен. \textit{Тема:} Обзор / книга.
\bibitem[44]{doi44}
C. Quinn, K. Rinne, T. O'Donnell, M. Duffy, C.O. Mathuna. \textit{A review of planar magnetic techniques and technologies}. .
\url{https://doi.org/10.1109/apec.2001.912514}.
\textit{Доступность:} Проверена библиографическая запись Crossref; полный текст не получен. \textit{Тема:} Обзор / книга.
\bibitem[45]{doi45}
E.C.W. de Jong, J.A. Ferreira, P. Bauer. \textit{Improving the thermal management of AC-DC converters using integration technologies}. .
\url{https://doi.org/10.1109/ias.2004.1348798}.
\textit{Доступность:} Проверена библиографическая запись Crossref; полный текст не получен. \textit{Тема:} Тепловой расчёт.
\bibitem[46]{doi46}
E. Waffenschmidt, J.A. Ferreira. \textit{Embedded passives integrated circuits for power converters}. .
\url{https://doi.org/10.1109/psec.2002.1023840}.
\textit{Доступность:} Проверена библиографическая запись Crossref; полный текст не получен. \textit{Тема:} Планарные компоненты и схемные модели.
\bibitem[47]{doi47}
K.D.T. Ngo, R.P. Alley, A.J. Yerman. \textit{Fabrication method for a winding assembly with a large number of planar layers}. 1993.
\url{https://doi.org/10.1109/63.208500}.
\textit{Доступность:} Проверена библиографическая запись Crossref; полный текст не получен. \textit{Тема:} Изготовление и интеграция.
\bibitem[48]{doi48}
M. Ludwig, M. Duffy, T. O'Donnell, P. McCloskey, S.C.O. Mathuna. \textit{PCB integrated inductors for low power DC/DC converter}. 2003.
\url{https://doi.org/10.1109/tpel.2003.813757}.
\textit{Доступность:} Проверена библиографическая запись Crossref; полный текст не получен. \textit{Тема:} Планарные компоненты и схемные модели.
\bibitem[49]{doi49}
E. Waffenschmidt, B. Ackermann, J.A. Ferreira. \textit{Design Method and Material Technologies for Passives in Printed Circuit Board Embedded Circuits}. 2005.
\url{https://doi.org/10.1109/tpel.2005.846530}.
\textit{Доступность:} Проверена библиографическая запись Crossref; полный текст не получен. \textit{Тема:} Изготовление и интеграция.
\bibitem[50]{doi50}
W.A. Roshen. \textit{Analysis of planar sandwich inductors by current images}. 1990.
\url{https://doi.org/10.1109/20.104901}.
\textit{Доступность:} Проверена библиографическая запись Crossref; полный текст не получен. \textit{Тема:} Планарные компоненты и схемные модели.
\bibitem[51]{doi51}
W.A. Roshen. \textit{Effect of finite thickness of magnetic substrate on planar inductors}. 1990.
\url{https://doi.org/10.1109/20.50553}.
\textit{Доступность:} Проверена библиографическая запись Crossref; полный текст не получен. \textit{Тема:} Планарные компоненты и схемные модели.
\bibitem[52]{doi52}
Y. Fukuda, T. Inoue, T. Mizoguchi, S. Yatabe, Y. Tachi. \textit{Planar inductor with ferrite layers for DC-DC converter}. 2003.
\url{https://doi.org/10.1109/tmag.2003.812708}.
\textit{Доступность:} Проверена библиографическая запись Crossref; полный текст не получен. \textit{Тема:} Планарные компоненты и схемные модели.
\bibitem[53]{doi53}
P. Wallmeier. \textit{Improved analytical modeling of conductive losses in gapped high-frequency inductors}. 2001.
\url{https://doi.org/10.1109/28.936395}.
\textit{Доступность:} Проверена библиографическая запись Crossref; полный текст не получен. \textit{Тема:} Потери и зазор.
\bibitem[54]{doi54}
E.C.W. de Jong, B.J.A. Ferreira, P. Bauer. \textit{Toward the Next Level of PCB Usage in Power Electronic Converters}. 2008.
\url{https://doi.org/10.1109/tpel.2008.2004276}.
\textit{Доступность:} Полный PDF получен из университетского репозитория; текст прочитан с ограничениями. \textit{Тема:} Изготовление и интеграция.
\bibitem[55]{doi55}
C. Buccella, C. Cecati, F. de Monte. \textit{A Coupled Electrothermal Model for Planar Transformer Temperature Distribution Computation}. 2008.
\url{https://doi.org/10.1109/tie.2008.2003102}.
\textit{Доступность:} Проверена библиографическая запись Crossref; полный текст не получен. \textit{Тема:} Тепловой расчёт.
\bibitem[56]{doi56}
A. Prasai, W.G. Odendaal. \textit{Fabrication and modeling of a planar magnetic structure with directly etched windings}. 2009.
\url{https://doi.org/10.1109/ecce.2009.5316202}.
\textit{Доступность:} Проверена библиографическая запись Crossref; полный текст не получен. \textit{Тема:} Изготовление и интеграция.
\bibitem[57]{doi57}
Roberto Prieto, Jesus A Oliver, Jose A Cobos, Mark Christini. \textit{Magnetic Component Model for Planar Structures Based on Transmission Lines}. 2010.
\url{https://doi.org/10.1109/tie.2009.2032434}.
\textit{Доступность:} Проверена библиографическая запись Crossref; полный текст не получен. \textit{Тема:} Планарные компоненты и схемные модели.
\bibitem[58]{doi58}
Alexander W Roesler, Joshua M Schare, S Jill Glass, Kevin G Ewsuk, George Slama, David Abel, Daryl Schofield. \textit{Planar LTCC Transformers for High-Voltage Flyback Converters}. 2010.
\url{https://doi.org/10.1109/tcapt.2009.2031872}.
\textit{Доступность:} Проверена библиографическая запись Crossref; полный текст не получен. \textit{Тема:} Планарные компоненты и схемные модели.
\bibitem[59]{doi59}
Chi Kwan Lee, Y. P. Su, S. Y. Ron Hui. \textit{Printed Spiral Winding Inductor With Wide Frequency Bandwidth}. 2011.
\url{https://doi.org/10.1109/tpel.2010.2076318}.
\textit{Доступность:} Проверена библиографическая запись Crossref; полный текст не получен. \textit{Тема:} Планарные компоненты и схемные модели.
\bibitem[60]{doi60}
Juergen Biela, Johann W. Kolar. \textit{Using Transformer Parasitics for Resonant Converters—A Review of the Calculation of the Stray Capacitance of Transformers}. 2008.
\url{https://doi.org/10.1109/tia.2007.912722}.
\textit{Доступность:} Проверена библиографическая запись Crossref; полный текст не получен. \textit{Тема:} Паразитные параметры и ЭМС.
\bibitem[61]{doi61}
Snezana Djuric, Goran Stojanovic, Mirjana Damnjanovic, Milan Radovanovic, Eric Laboure. \textit{Design, Modeling, and Analysis of a Compact Planar Transformer}. 2012.
\url{https://doi.org/10.1109/tmag.2012.2202642}.
\textit{Доступность:} Проверена библиографическая запись Crossref; полный текст не получен. \textit{Тема:} Планарные компоненты и схемные модели.
\bibitem[62]{doi62}
R. Chen, J.D. vanWyk, S. Wang, W.G. Odendaal. \textit{Improving the Characteristics of Integrated EMI Filters by Embedded Conductive Layers}. 2005.
\url{https://doi.org/10.1109/tpel.2005.846526}.
\textit{Доступность:} Проверена библиографическая запись Crossref; полный текст не получен. \textit{Тема:} Паразитные параметры и ЭМС.
\bibitem[63]{doi63}
Samuel R. Cove, Martin Ordonez, Federico Luchino, John E. Quaicoe. \textit{Applying Response Surface Methodology to Small Planar Transformer Winding Design}. 2013.
\url{https://doi.org/10.1109/tie.2012.2187416}.
\textit{Доступность:} Проверена библиографическая запись Crossref; полный текст не получен. \textit{Тема:} Планарные компоненты и схемные модели.
\bibitem[64]{doi64}
Junwei Lu, Francis Dawson. \textit{Characterizations of High Frequency Planar Transformer With a Novel Comb-Shaped Shield}. 2011.
\url{https://doi.org/10.1109/tmag.2011.2157664}.
\textit{Доступность:} Проверена библиографическая запись Crossref; полный текст не получен. \textit{Тема:} Паразитные параметры и ЭМС.
\bibitem[65]{doi65}
M.T. Quirke, J.J. Barrett, M. Hayes. \textit{Planar magnetic component technology-a review}. 1992.
\url{https://doi.org/10.1109/33.180055}.
\textit{Доступность:} Проверена библиографическая запись Crossref; полный текст не получен. \textit{Тема:} Обзор / книга.
\bibitem[66]{doi66}
Luca Dalessandro, Fabiana da Silveira Cavalcante, Johann W. Kolar. \textit{Self-Capacitance of High-Voltage Transformers}. 2007.
\url{https://doi.org/10.1109/tpel.2007.904252}.
\textit{Доступность:} Проверена библиографическая запись Crossref; полный текст не получен. \textit{Тема:} Паразитные параметры и ЭМС.
\bibitem[67]{doi67}
L. Heinemann. \textit{Modelling and design of high frequency planar transformers}. .
\url{https://doi.org/10.1109/pesc.1995.474888}.
\textit{Доступность:} Проверена библиографическая запись Crossref; полный текст не получен. \textit{Тема:} Планарные компоненты и схемные модели.
\bibitem[68]{doi68}
D. van der Linde, C.A.M. Boon, J.B. Klaassens. \textit{Design of a high-frequency planar power transformer in multilayer technology}. 1991.
\url{https://doi.org/10.1109/41.88907}.
\textit{Доступность:} Проверена библиографическая запись Crossref; полный текст не получен. \textit{Тема:} Изготовление и интеграция.
\bibitem[69]{doi69}
S.C. Tang, S.Y.R. Hui, H.S.-H. Chung. \textit{A low-profile power converter using printed-circuit board (PCB) power transformer with ferrite polymer composite}. 2001.
\url{https://doi.org/10.1109/63.931055}.
\textit{Доступность:} Проверена библиографическая запись Crossref; полный текст не получен. \textit{Тема:} Планарные компоненты и схемные модели.
\bibitem[70]{doi70}
K.D. Papastergiou, D.E. Macpherson. \textit{Air-gap effects in inductive energy transfer}. 2008.
\url{https://doi.org/10.1109/pesc.2008.4592594}.
\textit{Доступность:} Проверена библиографическая запись Crossref; полный текст не получен. \textit{Тема:} Потери и зазор.
\bibitem[71]{doi71}
J. Fletcher, B. Williams, M. Mahmoud. \textit{Airgap fringing flux reduction in inductors using open-circuit copper screens}. 2005.
\url{https://doi.org/10.1049/ip-epa:20045234}.
\textit{Доступность:} Проверена библиографическая запись Crossref; полный текст не получен. \textit{Тема:} Потери и зазор.
\bibitem[72]{doi72}
Jian Sun, V. Mehrotra. \textit{Orthogonal Winding Structures and Design for Planar Integrated Magnetics}. 2008.
\url{https://doi.org/10.1109/tie.2007.909754}.
\textit{Доступность:} Проверена библиографическая запись Crossref; полный текст не получен. \textit{Тема:} Планарные компоненты и схемные модели.
\bibitem[73]{doi73}
Ziwei Ouyang, Gökhan Sen, Ole C. Thomsen, Michael A. E. Andersen. \textit{Analysis and Design of Fully Integrated Planar Magnetics for Primary–Parallel Isolated Boost Converter}. 2013.
\url{https://doi.org/10.1109/tie.2012.2186777}.
\textit{Доступность:} Проверена библиографическая запись Crossref; полный текст не получен. \textit{Тема:} Планарные компоненты и схемные модели.
\bibitem[74]{doi74}
Xingkui Mao, Wei Chen, Yunxiu Li. \textit{Winding loss mechanism analysis and design for new structure high-frequency gapped inductor}. 2005.
\url{https://doi.org/10.1109/tmag.2005.854991}.
\textit{Доступность:} Проверена библиографическая запись Crossref; полный текст не получен. \textit{Тема:} Потери и зазор.
\bibitem[75]{doi75}
M. Albach, H. Rossmanith. \textit{The influence of air gap size and winding position on the proximity losses in high frequency transformers}. .
\url{https://doi.org/10.1109/pesc.2001.954329}.
\textit{Доступность:} Проверена библиографическая запись Crossref; полный текст не получен. \textit{Тема:} Потери и зазор.
\bibitem[76]{doi76}
W.A. Roshen. \textit{High-Frequency Fringing Fields Loss in Thick Rectangular and Round Wire Windings}. 2008.
\url{https://doi.org/10.1109/tmag.2008.2002302}.
\textit{Доступность:} Проверена библиографическая запись Crossref; полный текст не получен. \textit{Тема:} Потери и зазор.
\bibitem[77]{doi77}
W.A. Roshen. \textit{Fringing Field Formulas and Winding Loss Due to an Air Gap}. 2007.
\url{https://doi.org/10.1109/tmag.2007.898908}.
\textit{Доступность:} Проверена библиографическая запись Crossref; полный текст не получен. \textit{Тема:} Потери и зазор.
\bibitem[78]{ouyangreview}
Ziwei Ouyang, Michael A. E. Andersen. \textit{Overview of Planar Magnetic Technology—Fundamental Properties}. 2014.
\url{https://doi.org/10.1109/tpel.2013.2283263}.
\textit{Доступность:} Проверена библиографическая запись Crossref; полный текст не получен. \textit{Тема:} Обзор / книга.
\bibitem[79]{schaefer}
Jannik Schafer, Dominik Bortis, Johann W. Kolar. \textit{Optimal Design of Highly Efficient and Highly Compact PCB Winding Inductors}. 2018.
\url{https://doi.org/10.1109/compel.2018.8460166}.
\textit{Доступность:} Полный PDF получен; доступность не означает постраничного рецензирования. \textit{Тема:} Планарные компоненты и схемные модели.
\bibitem[80]{ewaldfoil}
Thomas Ewald, Jurgen Biela. \textit{Frequency-Dependent Inductance and Winding Loss Model for Gapped Foil Inductors}. 2022.
\url{https://doi.org/10.1109/tpel.2022.3169620}.
\textit{Доступность:} Полный PDF получен; доступность не означает постраничного рецензирования. \textit{Тема:} Потери и зазор.
\bibitem[81]{doi81}
Thomas Ewald, Jürgen Biela. \textit{Analytical Fringing Field Model for Efficient Winding Loss Calculation in Gapped Foil Inductors}. 2024.
\url{https://doi.org/10.1109/tpel.2024.3368145}.
\textit{Доступность:} Проверена библиографическая запись Crossref; полный текст не получен. \textit{Тема:} Потери и зазор.
\bibitem[82]{doi82}
Jannik Schafer, Dominik Bortis, Johann W. Kolar. \textit{Novel Highly Efficient/Compact Automotive PCB Winding Inductors Based on the Compensating Air-Gap Fringing Field Concept}. 2020.
\url{https://doi.org/10.1109/tpel.2020.2969295}.
\textit{Доступность:} Проверена библиографическая запись Crossref; полный текст не получен. \textit{Тема:} Потери и зазор.
\bibitem[83]{doi83}
Shenli Zou, Chanaka Singhabahu, Jianfei Chen, Alireza Khaligh. \textit{A comprehensive design approach for a three‐winding planar transformer}. 2022.
\url{https://doi.org/10.1049/pel2.12261}.
\textit{Доступность:} Проверена библиографическая запись Crossref; полный текст не получен. \textit{Тема:} Планарные компоненты и схемные модели.
\bibitem[84]{doi84}
Colonel Wm. T. McLyman. \textit{Transformer and Inductor Design Handbook}. 2017.
\url{https://doi.org/10.1201/b10865}.
\textit{Доступность:} Проверена библиографическая запись Crossref; полный текст не получен. \textit{Тема:} Обзор / книга.
\bibitem[85]{doi85}
Alexandre Msellati, Matthew Jahnes, Matthias Preindl. \textit{I-core PCB Planar Inductor Design for High Frequency and High Power Converters}. 2022.
\url{https://doi.org/10.1109/itec53557.2022.9813880}.
\textit{Доступность:} Проверена библиографическая запись Crossref; полный текст не получен. \textit{Тема:} Планарные компоненты и схемные модели.
\bibitem[86]{doi86}
Ruiwen Zhang, Daming Zhang, Rukmi Dutta. \textit{Study on PCB Based Litz Wire Applications for Air-Core Inductor and Planar Transformer}. 2019.
\url{https://doi.org/10.1109/icpes47639.2019.9105549}.
\textit{Доступность:} Проверена библиографическая запись Crossref; полный текст не получен. \textit{Тема:} Планарные компоненты и схемные модели.
\bibitem[87]{zhaocap}
Hongbo Zhao et al.. \textit{Parasitic Capacitance Modeling of Copper-Foiled Medium-Voltage Filter Inductors Considering Fringe Electrical Field}. 2021.
\url{https://doi.org/10.1109/TPEL.2020.3048226}.
\textit{Доступность:} Полный PDF получен; доступность не означает постраничного рецензирования. \textit{Тема:} Паразитные параметры и ЭМС.
\bibitem[88]{local4}
\textit{В. П. Обрусник. Магнитные элементы электронных устройств. Учебное пособие. 2018}. .
Локальный учебник из базы знаний проекта; путь приведён в машиночитаемом реестре.
\textit{Доступность:} Полный PDF получен; доступность не означает постраничного рецензирования. \textit{Тема:} Учебник: общий магнитный расчёт.
\bibitem[89]{local5}
\textit{Н. С. Легостаев. Магнитные элементы электронных устройств. Учебное пособие. 2014}. .
Локальный учебник из базы знаний проекта; путь приведён в машиночитаемом реестре.
\textit{Доступность:} Полный PDF получен; доступность не означает постраничного рецензирования. \textit{Тема:} Учебник: общий магнитный расчёт.
\end{thebibliography}
\end{document}
